% ROZDZIAL 13 — Kontenery STL i algorytmy w C++20.
% Fragment do zbiorczego skryptu (bez \documentclass i \begin{document}).
% Wymagania: listings, amsmath, amssymb, booktabs, tabularx, array, hyperref.
% Kompilator kodu: C++20; nowe elementy standardu sa oznaczane w tekscie.
% Dolaczenie: % ROZDZIAL 13 — Kontenery STL i algorytmy w C++20.
% Fragment do zbiorczego skryptu (bez \documentclass i \begin{document}).
% Wymagania: listings, amsmath, amssymb, booktabs, tabularx, array, hyperref.
% Kompilator kodu: C++20; nowe elementy standardu sa oznaczane w tekscie.
% Dolaczenie: % ROZDZIAL 13 — Kontenery STL i algorytmy w C++20.
% Fragment do zbiorczego skryptu (bez \documentclass i \begin{document}).
% Wymagania: listings, amsmath, amssymb, booktabs, tabularx, array, hyperref.
% Kompilator kodu: C++20; nowe elementy standardu sa oznaczane w tekscie.
% Dolaczenie: % ROZDZIAL 13 — Kontenery STL i algorytmy w C++20.
% Fragment do zbiorczego skryptu (bez \documentclass i \begin{document}).
% Wymagania: listings, amsmath, amssymb, booktabs, tabularx, array, hyperref.
% Kompilator kodu: C++20; nowe elementy standardu sa oznaczane w tekscie.
% Dolaczenie: \input{rozdzial_13_kontenery_stl_algorytmy_cpp.tex}

\chapter{Kontenery STL i algorytmy}
\label{chap:cpp-stl-containers}

\section{Cel i zakres rozdziału}

Kontener przechowuje grupę elementów i pozwala je odnajdywać lub
zmieniać. Najprostsza mapa wyboru wygląda tak: \texttt{vector} jest
listą z dostępem przez indeks, \texttt{map} słownikiem klucz--wartość,
a \texttt{set} zbiorem wartości bez powtórzeń. Jeśli kolejność kluczy
nie ma znaczenia, istnieją także haszujące warianty
\texttt{unordered\_map} i \texttt{unordered\_set}.

Biblioteka standardowa C++ udostępnia gotowe struktury danych, których wybór ma
istotne konsekwencje dla złożoności operacji, lokalności pamięci, bezpieczeństwa
wyjątkowego i poprawności iteratorów. Samo rozróżnienie między
\texttt{std::vector}, \texttt{std::list} i \texttt{std::map} nie wystarcza:
trzeba rozumieć również \emph{własność elementów}, \emph{unieważnianie
iteratorów}, wymagania dotyczące porządku i haszowania oraz relację między
kontenerami a algorytmami. W dalszej części domyślnie używamy C++20.

\paragraph{Główna zasada projektowa.}
Najpierw ustalamy rzeczywiste operacje: czy potrzebujemy dostępu przez
indeks, wyszukiwania po kluczu, zachowania kolejności, stabilnych adresów,
czy częstych wstawek i usunięć. Dopiero potem wybieramy kontener. Prosty
\texttt{std::vector} pozostaje często najlepszym rozwiązaniem, nawet gdy
jego teoretyczna złożoność niektórych operacji jest gorsza, ponieważ dane
leżą w sąsiadujących komórkach pamięci.

\section{Klasyfikacja kontenerów i podstawowe pojęcia}
\label{sec:cpp13-taxonomy}

Biblioteka standardowa rozróżnia m.in. kontenery sekwencyjne
(\texttt{array}, \texttt{vector}, \texttt{deque}, \texttt{list},
\texttt{forward\_list}), uporządkowane asocjacyjne (\texttt{map},
\texttt{set} i ich warianty \texttt{multi}) oraz nieuporządkowane
asocjacyjne (\texttt{unordered\_map}, \texttt{unordered\_set} i warianty
\texttt{multi}). \texttt{stack}, \texttt{queue} i \texttt{priority\_queue}
są \emph{adapterami kontenerów}: oferują inny interfejs dostępu do
wewnętrznej struktury danych. \texttt{std::string} jest wyspecjalizowaną
strukturą sekwencyjną, ale semantykę tekstu omawiamy oddzielnie.

\begin{center}
\small
\begin{tabularx}{\textwidth}{@{}lXXX@{}}
\toprule
\textbf{Kontener} & \textbf{Dostęp} & \textbf{Wstawienie / usunięcie} & \textbf{Typowe zastosowanie} \\
\midrule
\texttt{vector} & Indeks $O(1)$ & Koniec: zamortyzowane $O(1)$; środek: $O(n)$ & Dane zwarte, iteracja, obliczenia \\
\texttt{deque} & Indeks $O(1)$ & Oba końce: $O(1)$; środek: $O(n)$ & Kolejka dwustronna \\
\texttt{list} & Iteracja $O(n)$ & $O(1)$ przy znanym iteratorze & Stabilne węzły, \texttt{splice} \\
\texttt{forward\_list} & Iteracja $O(n)$ & $O(1)$ za poprzednikiem & Lista jednokierunkowa \\
\texttt{map} / \texttt{set} & Szukanie $O(\log n)$ & $O(\log n)$ & Klucze w ustalonej kolejności \\
\texttt{unordered\_map} & Średnio $O(1)$ & Średnio $O(1)$ & Słownik, szybkie wyszukiwanie \\
\bottomrule
\end{tabularx}
\end{center}

W tabeli $n$ oznacza liczbę elementów, a koszty zależą od wymagań
poszczególnych operacji. \texttt{unordered\_map} może w najgorszym
przypadku wykonywać operacje w czasie $O(n)$ (np. przy wielu kolizjach).
Sformułowanie ,,\emph{zamortyzowane} $O(1)$'' oznacza, że pojedyncza
operacja może być kosztowna, ale średni koszt w długiej sekwencji operacji
jest ograniczony stałą.

\paragraph{Wartości, wskaźniki, referencje.}
Kontenery zwykle przechowują \emph{własne obiekty elementów}.
\texttt{std::vector<T>} nie jest kontenerem referencji do obiektów
\texttt{T}; może za to przechowywać np. \texttt{std::reference\_wrapper<T>},
\texttt{T*} lub \texttt{std::unique\_ptr<T>}. Wybór wpływa na własność
obiektów, kopiowanie, czas życia i poprawność dostępu.

\section{\texttt{std::vector}: domyślny kontener sekwencyjny}

\begin{center}\small
\begin{tabularx}{\textwidth}{@{}p{0.34\textwidth}X@{}}
\toprule
Operacja & Za co odpowiada \\
\midrule
\code{size()}, \code{empty()} & Podają liczbę elementów i informację, czy kontener jest pusty. \\
\code{capacity()} & Podaje liczbę elementów, które zmieszczą się bez ponownej alokacji. \\
\code{reserve(n)} & Rezerwuje miejsce na co najmniej \code{n} elementów; nie zmienia \code{size()}. \\
\code{resize(n)} & Zmienia liczbę elementów: tworzy nowe albo niszczy nadmiarowe. \\
\code{push_back(x)} & Dodaje istniejącą wartość na koniec przez kopiowanie lub przeniesienie. \\
\code{emplace_back(args...)} & Konstruuje nowy element na końcu z podanych argumentów. \\
\code{erase(pos)} & Usuwa element; kolejne elementy przesuwają się, a iteratory od miejsca usunięcia tracą ważność. \\
\bottomrule
\end{tabularx}
\end{center}
\label{sec:cpp13-vector}

\texttt{std::vector<T>} przechowuje elementy \emph{ciągle} w pamięci,
co pozwala przekazywać je do interfejsów oczekujących tablicy (np. przez
\texttt{data()} lub \texttt{std::span}). Obiekt wektora zarządza własną
pamięcią i niszczy elementy po zakończeniu swojego czasu życia.

\begin{lstlisting}[caption={Podstawowe operacje na vector (C++11); kompletny program.}]
#include <cassert>
#include <vector>

int main() {
    std::vector<int> v{4, 8, 15};
    v.push_back(16);           // kopiowanie lub przenoszenie argumentu
    v.emplace_back(23);        // konstruowanie nowego elementu

    assert(v.size() == 5);
    assert(v[0] == 4);
    assert(v.at(4) == 23);
    assert(v.front() == 4 && v.back() == 23);

    v.pop_back();
    assert(v.size() == 4);
    v.clear();
    assert(v.empty());
}
\end{lstlisting}

\texttt{operator[]} \emph{nie sprawdza} granic: użycie indeksu spoza zakresu
prowadzi do niezdefiniowanego zachowania. Metoda \texttt{at()} sprawdza
zakres i zgłasza \texttt{std::out\_of\_range}. Metody \texttt{front()},
\texttt{back()} i \texttt{pop\_back()} wymagają niepustego kontenera.

\subsection{\texttt{size}, \texttt{capacity}, \texttt{reserve} i \texttt{resize}}

\texttt{size()} oznacza liczbę \emph{istniejących} elementów, a
\texttt{capacity()} liczbę elementów, dla których przydzielono miejsce
bez konieczności realokacji. \texttt{reserve(k)} zmienia rezerwację
pamięci, nie tworzy elementów; \texttt{resize(k)} zmienia ich liczbę.

\begin{lstlisting}[caption={Rezerwacja pamięci nie zwiększa rozmiaru (C++11); kompletny program.}]
#include <cassert>
#include <vector>

int main() {
    std::vector<int> data;
    data.reserve(100);
    assert(data.size() == 0);
    assert(data.capacity() >= 100);

    data.push_back(7);
    data.resize(4);          // dopisuje trzy elementy o wartosci 0
    assert(data.size() == 4);
    assert(data[0] == 7 && data[3] == 0);
}
\end{lstlisting}

Nie należy wywoływać \texttt{reserve(size()+1)} przed każdym
\texttt{push\_back()}: wymuszenie niemal liniowego wzrostu pojemności
może zniweczyć politykę geometrycznego powiększania bufora i pogorszyć
wydajność. \texttt{shrink\_to\_fit()} jest tylko żądaniem, a nie gwarancją
zmniejszenia pojemności.

\subsection{Realokacja i unieważnianie iteratorów}

Po realokacji \texttt{vector} zmienia adresy wszystkich elementów:
tracą ważność iteratory, referencje i wskaźniki do elementów, również
poprzednia wartość \texttt{data()}. Bez realokacji \texttt{push\_back()}
unieważnia iterator \texttt{end()}, lecz nie unieważnia iteratorów i
referencji do wcześniej istniejących elementów. Wstawienie w środku
bez realokacji unieważnia iteratory i referencje do pozycji wstawienia
oraz wszystkich późniejszych pozycji (włącznie z dawnym \texttt{end()}).

\begin{lstlisting}[caption={Bezpieczne ponowne pobranie iteratora po wstawieniu; kompletny program.}]
#include <cassert>
#include <vector>

int main() {
    std::vector<int> v{1, 3, 4};
    auto pos = v.begin() + 1;
    pos = v.insert(pos, 2);   // przyjmuje zwrocony iterator
    assert(*pos == 2);
    assert((v == std::vector<int>{1, 2, 3, 4}));
}
\end{lstlisting}

\paragraph{Przenoszenie elementów i bezpieczeństwo wyjątkowe.}
Podczas realokacji \texttt{vector} może przenosić lub kopiować elementy.
Dla pewnych typów i operacji warunki silnej gwarancji wyjątkowej zależą
od tego, czy konstruktor przenoszący jest \texttt{noexcept} oraz czy
dostępne jest kopiowanie. Zagadnienie łączy się bezpośrednio z semantyką
przenoszenia omówioną w rozdziale 6.

\subsection{\texttt{push\_back} kontra \texttt{emplace\_back}}

\texttt{emplace\_back(args...)} przekazuje argumenty bezpośrednio do
konstruktora nowego elementu. Nie oznacza to jednak automatycznie
,,zawsze szybciej'': jeśli mamy już obiekt, \texttt{push\_back(std::move(x))}
bywa równie naturalne; przy realokacji nadal konieczne może być
przeniesienie innych elementów. Obie funkcje mogą zgłaszać wyjątki.

\begin{lstlisting}[caption={Konstruowanie elementu z argumentow (C++17); kompletny program.}]
#include <cassert>
#include <string>
#include <utility>
#include <vector>

struct Person {
    std::string name;
    int age;
    Person(std::string n, int a) : name(std::move(n)), age(a) {}
};

int main() {
    std::vector<Person> people;
    people.emplace_back("Ada", 31);
    people.push_back(Person{"Alan", 42});
    assert(people[1].age == 42);
}
\end{lstlisting}

\section{\texttt{std::array}, \texttt{std::deque} i listy}
\label{sec:cpp13-sequential}

\subsection{\texttt{std::array<T,N>}: rozmiar znany w kompilacji}

\texttt{std::array} przechowuje dokładnie $N$ elementów. Nie dokonuje
osobnej dynamicznej alokacji dla swojej zawartości; pamięć elementów
stanowi część samego obiektu \texttt{array}. Można używać iteratorów,
\texttt{size()}, \texttt{data()} i algorytmów, analogicznie do wektora.

\begin{lstlisting}[caption={std::array i algorytmy standardowe (C++17); kompletny program.}]
#include <algorithm>
#include <array>
#include <cassert>
#include <tuple>

int main() {
    std::array<int, 4> x{5, 1, 4, 2};
    std::sort(x.begin(), x.end());
    assert((x == std::array<int, 4>{1, 2, 4, 5}));
    static_assert(std::tuple_size_v<decltype(x)> == 4);
}
\end{lstlisting}

\subsection{\texttt{std::deque}: dwa tanie końce}

\texttt{deque} udostępnia dostęp indeksowany w czasie $O(1)$, lecz
zwykle nie jest jednym ciągłym blokiem pamięci. Wstawianie i usuwanie
na początku i końcu ma stałą złożoność; w przeciwieństwie do
\texttt{vector} ma operacje \texttt{push\_front()} i \texttt{pop\_front()}.
\texttt{deque} nie udostępnia \texttt{data()} dla wszystkich elementów.

\begin{lstlisting}[caption={Kolejka dwustronna (C++11); kompletny program.}]
#include <cassert>
#include <deque>

int main() {
    std::deque<int> q{2, 3};
    q.push_front(1);
    q.push_back(4);
    assert(q[0] == 1 && q[3] == 4);
    q.pop_front();
    assert(q.front() == 2);
}
\end{lstlisting}

Wstawienie na końcu \texttt{deque} może unieważnić \emph{iteratory},
ale zachowuje referencje do istniejących elementów. Szczegółowe reguły
różnią się dla wstawień i usunięć w środku oraz na końcach; nie należy
mechanicznie przenosić reguł \texttt{vector} na \texttt{deque}.

\subsection{\texttt{std::list} i \texttt{std::forward\_list}}

\texttt{list} jest listą dwukierunkową. Wstawianie i usuwanie przy
znanym iteratorze odbywa się w stałym czasie; iteratory i referencje
do pozostałych elementów pozostają ważne. Wyszukanie elementu przez
wartość wymaga jednak przejścia listy $O(n)$, podobnie jak uzyskanie
elementu o zadanym indeksie. W przeciwieństwie do \texttt{vector},
\texttt{list} nie oferuje operatora \texttt{[]} ani iteratorów dostępu
swobodnego. Do sortowania używamy \texttt{list::sort()}, a nie
\texttt{std::sort()}.

\begin{lstlisting}[caption={Przeniesienie wezla przez list::splice (C++11); kompletny program.}]
#include <cassert>
#include <iterator>
#include <list>

int main() {
    std::list<int> source{1, 2, 3};
    std::list<int> target{10};
    auto it = std::next(source.begin()); // wskazuje 2

    target.splice(target.end(), source, it);
    assert(*it == 2); // iterator pozostaje wazny
    assert(source.size() == 2 && target.size() == 2);
    target.sort();
    assert(target.front() == 2);
}
\end{lstlisting}

\texttt{forward\_list} przechowuje węzły listy jednokierunkowej;
wstawianie i usuwanie wykonuje \emph{za} wskazanym poprzednikiem
(\texttt{insert\_after}, \texttt{erase\_after}), a \texttt{before\_begin()}
udostępnia pozycję przed pierwszym elementem. Kontener nie posiada
metody \texttt{size()} (rozmiar można policzyć liniowo przez
\texttt{std::distance}).

\paragraph{Kiedy lista jest rzeczywiście potrzebna?}
Lista bywa przydatna, gdy musimy często \emph{przenosić węzły bez
kopiowania obiektów} i wymagamy stabilnych iteratorów. Dla częstej
iteracji i obliczeń numerycznych \texttt{vector} z reguły zapewnia
lepszą lokalność pamięci. Wybór należy jednak potwierdzić pomiarem.

\section{Kontenery uporządkowane: \texttt{set} i \texttt{map}}
\label{sec:cpp13-ordered}

\texttt{std::set<Key>} przechowuje unikalne klucze, a
\texttt{std::map<Key,T>} unikalne pary klucz--wartość. Elementy są
iterowane w kolejności zdefiniowanej przez komparator; wyszukiwanie,
wstawianie i usuwanie są logarytmiczne. Warianty
\texttt{multiset}/\texttt{multimap} dopuszczają powtarzające się klucze.

\subsection{\texttt{map::operator[]}, \texttt{at} i \texttt{find}}

Odwołanie \texttt{m[key]} do nieistniejącego klucza \emph{wstawia} nowy
element z wartością inicjalizowaną domyślnie; może zatem modyfikować
strukturę danych. \texttt{at()} zgłasza wyjątek, jeżeli klucz nie
istnieje. \texttt{find()} zwraca iterator, a \texttt{contains()} (C++20)
zwraca wartość logiczną bez wstawiania.

\begin{lstlisting}[caption={Liczenie czestosci slow z map (C++20); kompletny program.}]
#include <cassert>
#include <map>
#include <string>
#include <vector>

int main() {
    std::vector<std::string> words{"kot", "pies", "kot"};
    std::map<std::string, int> count;
    for (const auto& word : words) ++count[word];

    assert(count.at("kot") == 2);
    assert(count.contains("pies"));
    assert(!count.contains("lis"));
    assert(count.find("lis") == count.end());
}
\end{lstlisting}

\subsection{\texttt{insert}, \texttt{try\_emplace}, \texttt{insert\_or\_assign}}

\texttt{insert} podejmuje próbę dodania gotowej pary.
\texttt{try\_emplace(key,args...)} (C++17) konstruuje wartość
w miejscu z przekazanych argumentów tylko wtedy, gdy klucz nie występuje.
Ważne zastrzeżenie: \emph{wyrażenia będące argumentami wywołania}
są obliczane jeszcze przed wejściem do \texttt{try\_emplace}; nie
należy więc podawać niepotrzebnie kosztownego wyrażenia tworzącego
tymczasowy obiekt. \texttt{insert\_or\_assign} dodaje lub zastępuje wartość.

\begin{lstlisting}[caption={Rozne operacje modyfikacji map (C++17); kompletny program.}]
#include <cassert>
#include <map>
#include <memory>
#include <string>

struct Heavy {
    static inline int constructions = 0; // C++17
    int value;
    explicit Heavy(int x) : value(x) { ++constructions; }
};

int main() {
    std::map<std::string, Heavy> m;
    auto [it, added] = m.try_emplace("a", 3);
    assert(added && it->second.value == 3);

    auto [again, added_again] = m.try_emplace("a", 4);
    assert(!added_again && again->second.value == 3);
    assert(Heavy::constructions == 1);

    m.insert_or_assign("a", Heavy{9});
    assert(m.at("a").value == 9);

    std::map<int, std::unique_ptr<int>> owned;
    owned.try_emplace(1, std::make_unique<int>(10));
    assert(*owned.at(1) == 10);
}
\end{lstlisting}

\subsection{Własny komparator i ścisły słaby porządek}

Porównanie musi spełniać warunek \emph{strict weak ordering}; w
szczególności komparator powinien być irrefleksyjny
($\neg\operatorname{comp}(x,x)$) i przechodni. Jeżeli dwie wartości
$ a,b $ nie spełniają ani \texttt{comp(a,b)}, ani \texttt{comp(b,a)},
kontener uznaje je za równoważne jako klucze --- nawet gdy
\texttt{operator==} zwraca \texttt{false}. Wymagania tego porządku
są istotne również dla \texttt{sort} i \texttt{lower\_bound}.
Na przykład komparator ignorujący wielkość liter uzna \texttt{"ABC"}
i \texttt{"abc"} za ten sam klucz: druga próba wstawienia do
\texttt{map} nie utworzy nowego elementu.

\begin{lstlisting}[caption={Mapowanie bez rozrozniania wielkosci liter: prosty komparator; kompletny program.}]
#include <cassert>
#include <cctype>
#include <map>
#include <string>

struct CaseInsensitiveLess {
    bool operator()(const std::string& a,
                    const std::string& b) const {
        auto ai = a.begin();
        auto bi = b.begin();
        while (ai != a.end() && bi != b.end()) {
            auto x = static_cast<unsigned char>(*ai++);
            auto y = static_cast<unsigned char>(*bi++);
            int lx = std::tolower(x);
            int ly = std::tolower(y);
            if (lx != ly) return lx < ly;
        }
        return a.size() < b.size();
    }
};

int main() {
    std::map<std::string, int, CaseInsensitiveLess> m;
    m.emplace("ABC", 1);
    m.emplace("abc", 2); // rownowazny klucz
    assert(m.size() == 1 && m.begin()->second == 1);
}
\end{lstlisting}

Przykład porównuje bajty w podstawowym środowisku znakowym;
\texttt{std::tolower} z bieżącą lokalizacją C nie realizuje ogólnego
Unicode case folding. Do pełnego porównywania napisów w Unicode trzeba
zastosować osobne, świadomie wybrane reguły normalizacji i kolacji.

\subsection{Wyszukiwanie heterogeniczne}

Jeżeli komparator mapy jest \emph{transparentny} (np. \texttt{std::less<>}),
możemy szukać w mapie o kluczach \texttt{std::string} za pomocą
\texttt{std::string\_view}, bez tworzenia tymczasowego ciągu znaków.
Heterogeniczne przeciążenia \texttt{find} są dostępne
od C++14, a \texttt{contains} od C++20.

\begin{lstlisting}[caption={Heterogeniczne wyszukiwanie w map (C++20); kompletny program.}]
#include <cassert>
#include <functional>
#include <map>
#include <string>
#include <string_view>

int main() {
    std::map<std::string, int, std::less<>> m{{"alpha", 3}};
    std::string_view key = "alpha";
    auto it = m.find(key);
    assert(it != m.end() && it->second == 3);
    assert(m.contains(key));
}
\end{lstlisting}

\subsection{Zmiana klucza przez uchwyt węzła (C++17)}

Kluczy w \texttt{map} nie wolno modyfikować bezpośrednio przez iterator.
Uchwyt \texttt{node\_handle} uzyskany przez \texttt{extract()} pozwala
wyjąć węzeł, zmienić jego klucz i wstawić z powrotem. Trzeba pamiętać,
że ponowne wstawienie może się nie udać z powodu kolizji kluczy.

\begin{lstlisting}[caption={Zmiana klucza przez extract (C++17); kompletny program.}]
#include <cassert>
#include <map>
#include <string>
#include <utility>

int main() {
    std::map<std::string, int> m{{"old", 7}};
    auto node = m.extract("old");
    assert(!node.empty());
    node.key() = "new";
    auto result = m.insert(std::move(node));
    assert(result.inserted && m.at("new") == 7);
}
\end{lstlisting}

\section{Kontenery haszujące: \texttt{unordered\_map} i \texttt{unordered\_set}}
\label{sec:cpp13-unordered}

Kontener \texttt{unordered\_map<Key,T,Hash,Eq>} rozmieszcza elementy
w kubełkach. Funkcja \texttt{Hash} przekształca klucz w wartość typu
\texttt{std::size\_t}; \texttt{Eq} rozstrzyga równość kluczy.
Wymagany jest warunek zgodności:
\[
\operatorname{Eq}(a,b) \Longrightarrow
\operatorname{Hash}(a)=\operatorname{Hash}(b).
\]
Warunek odwrotny nie zachodzi: kolizje haszy są dozwolone i
rozstrzygane przez porównanie kluczy. Iteracja po kontenerze
nie ma stabilnego porządku i nie należy z niego korzystać jako z
interfejsu logicznego programu.
Można wyobrazić sobie hasz jako wskazówkę, w którym kubełku szukać.
Jeśli dwa różne klucze trafią do tego samego kubełka, \texttt{Eq}
sprawdza, czy są naprawdę równe. Stała w poniższym \texttt{PointHash}
służy jedynie do wymieszania haszy obu współrzędnych.

\begin{lstlisting}[caption={Wlasny hash i porownanie kluczy (C++20); kompletny program.}]
#include <cassert>
#include <cstddef>
#include <functional>
#include <unordered_map>

struct Point {
    int x, y;
    friend bool operator==(const Point&, const Point&) = default;
};

struct PointHash {
    std::size_t operator()(const Point& p) const noexcept {
        std::size_t h1 = std::hash<int>{}(p.x);
        std::size_t h2 = std::hash<int>{}(p.y);
        return h1 ^ (h2 + std::size_t{0x9e3779b9}
                        + (h1 << 6) + (h1 >> 2));
    }
};

int main() {
    std::unordered_map<Point, int, PointHash> m;
    m[{2, 5}] = 42;
    assert(m.at(Point{2, 5}) == 42);
}
\end{lstlisting}

Ta funkcja mieszająca służy demonstracji składni, a nie zastosowaniom
kryptograficznym. Dla publicznie sterowanych kluczy przy projektowaniu
systemów wystawionych na ataki kolizyjne należy rozważyć również
odporność na złośliwe dane wejściowe.

\subsection{\texttt{load\_factor}, \texttt{reserve} i \texttt{rehash}}

\texttt{load\_factor()} to stosunek liczby elementów do liczby kubełków.
\texttt{max\_load\_factor()} podaje lub ustawia pożądany maksymalny
współczynnik. \texttt{reserve(n)} żąda przygotowania na $n$ elementów,
a \texttt{rehash(k)} na przynajmniej $k$ kubełków, z uwzględnieniem
limitów wynikających z współczynnika wypełnienia.

\begin{lstlisting}[caption={Rezerwowanie kubelek przed licznymi wstawkami (C++20); kompletny program.}]
#include <cassert>
#include <unordered_set>

int main() {
    std::unordered_set<int> seen;
    seen.reserve(1000);
    for (int i = 0; i < 1000; ++i) seen.insert(i);
    assert(seen.size() == 1000);
    assert(seen.contains(17));
}
\end{lstlisting}

Rehashing unieważnia iteratory, lecz \emph{nie unieważnia} wskaźników
i referencji do elementów. Wstawienie może unieważnić iteratory, jeśli
wywoła rehash. Usunięcie unieważnia iteratory i referencje do usuniętych
elementów. Nie wolno modyfikować kluczy w sposób zaburzający ich hash
lub równość podczas przebywania elementów w kontenerze.

\section{Adaptery: \texttt{stack}, \texttt{queue}, \texttt{priority\_queue}}
\label{sec:cpp13-adaptors}

\texttt{stack} udostępnia schemat LIFO, \texttt{queue} --- FIFO,
a \texttt{priority\_queue} --- dostęp do elementu o najwyższym
priorytecie według komparatora (domyślnie do największego).
Adaptery zwykle nie zapewniają bezpośredniej iteracji po elementach.

\begin{lstlisting}[caption={Trzy adaptery standardowe (C++11); kompletny program.}]
#include <cassert>
#include <functional>
#include <queue>
#include <stack>
#include <vector>

int main() {
    std::stack<int> s;
    s.push(1); s.push(2);
    assert(s.top() == 2);
    s.pop();

    std::queue<int> q;
    q.push(1); q.push(2);
    assert(q.front() == 1);
    q.pop();

    std::priority_queue<int, std::vector<int>, std::greater<int>> pq;
    pq.push(7); pq.push(2); pq.push(5);
    assert(pq.top() == 2); // kopiec minimum
}
\end{lstlisting}

\texttt{priority\_queue} zwykle opiera się na kopcu: \texttt{top()} ma
koszt $O(1)$, a \texttt{push()} i \texttt{pop()} $O(\log n)$. Nie jest
to kolejka FIFO ani struktura z aktualizacją priorytetu dowolnego elementu
przez gotowy publiczny interfejs.

\section{Iteratory, stabilność adresów i unieważnianie}
\label{sec:cpp13-invalidation}

Iterator nie jest ogólnie tożsamy ze wskaźnikiem. Może przechowywać
stan niezbędny do poruszania się po drzewie lub kubełkach. Po operacji
modyfikującej trzeba wiedzieć, czy iterator, referencja i wskaźnik do
elementu pozostają ważne. Typowe reguły podsumowano w tabeli.

\begin{center}
\footnotesize
\begin{tabularx}{\textwidth}{@{}lXX@{}}
\toprule
\textbf{Kontener} & \textbf{Wstawienie} & \textbf{Usunięcie} \\
\midrule
\texttt{vector} & Przy realokacji: wszystkie iteratory i referencje; bez realokacji: pozycja wstawienia i dalej & Pozycja usunięcia i dalej (także dawny \texttt{end()}) \\
\texttt{deque} & Na końcach: iteratory tracą ważność, referencje do elementów zwykle pozostają; w środku: także referencje & Zależy od pozycji; usunięcie w środku unieważnia wszystkie iteratory i referencje \\
\texttt{list} & Istniejące iteratory i referencje zachowane & Tylko usunięte elementy \\
\texttt{map}, \texttt{set} & Istniejące iteratory i referencje zachowane & Tylko usunięte elementy \\
\texttt{unordered\_*} & Iteratory mogą utracić ważność przy rehash; referencje zachowane & Tylko usunięte elementy \\
\bottomrule
\end{tabularx}
\end{center}

Dla \texttt{deque} szczególnie ostrożnie trzeba traktować
\texttt{end()}, którego ważność zależy od rodzaju usunięcia.
W tabeli podano reguły najważniejsze dla praktyki, a nie kompletną
specyfikację każdej przeciążonej operacji. Po operacjach zmieniających
strukturę bezpieczną praktyką jest korzystanie z iteratora zwróconego
przez metodę (np. \texttt{erase}, \texttt{insert}).

\subsection{Bezpieczne usuwanie podczas iteracji}

\begin{lstlisting}[caption={Usuwanie elementow map podczas iteracji (C++20); kompletny program.}]
#include <cassert>
#include <map>

int main() {
    std::map<int, int> m{{1, 2}, {2, 3}, {3, 4}};
    for (auto it = m.begin(); it != m.end();) {
        if (it->second % 2 == 0) it = m.erase(it);
        else ++it;
    }
    assert(m.size() == 1 && m.contains(2));
}
\end{lstlisting}

\paragraph{Szczególny przypadek \texttt{vector<bool>}.}
\texttt{std::vector<bool>} jest specjalizacją oszczędzającą pamięć,
która może reprezentować wartości w postaci bitów. Jej \texttt{operator[]}
nie musi zwracać \texttt{bool\&}, tylko specjalny obiekt proxy. Nie
należy oczekiwać normalnej tablicy obiektów \texttt{bool}, ciągłego
przechowywania pojedynczych \texttt{bool} ani bezpiecznej modyfikacji
odrębnych bitów z różnych wątków bez synchronizacji.

\section{Algorytmy standardowe i wymagania wobec iteratorów}

\begin{center}\small
\begin{tabularx}{\textwidth}{@{}p{0.34\textwidth}X@{}}
\toprule
Algorytm & Działanie i warunki \\
\midrule
\code{std::find} & Szuka pierwszej wartości równej wzorcowi; przy braku zwraca iterator końcowy. \\
\code{std::transform} & Stosuje funkcję do zakresu i zapisuje wyniki przez iterator wyjściowy. \\
\code{std::sort} & Sortuje zakres z iteratorami dostępu swobodnego; nie zachowuje kolejności elementów równoważnych. \\
\code{std::stable_sort} & Sortuje, zachowując wzajemną kolejność elementów równoważnych. \\
\code{std::lower_bound} & Szuka pierwszego elementu nie mniejszego od wzorca; wymaga odpowiednio uporządkowanego zakresu. \\
\code{std::erase_if} & Usuwa z kontenera elementy spełniające predykat; dostępne od C++20 dla wspieranych kontenerów. \\
\bottomrule
\end{tabularx}
\end{center}
\label{sec:cpp13-algorithms}

Algorytmy oddzielają \emph{operację} od \emph{reprezentacji danych}.
Wersje klasyczne przyjmują iteratory, a w C++20 wiele algorytmów
w \texttt{std::ranges} można wywołać bezpośrednio na zakresie. Szczegóły
kategorii iteratorów i widoków omawiamy w rozdziale~\ref{chap:ranges-views-cpp}; tutaj skupiamy
się na doborze algorytmu i jego kontrakcie.

\subsection{Wyszukiwanie, transformacje i agregacje}

Do częstych operacji należą \texttt{find}, \texttt{find\_if},
\texttt{count\_if}, \texttt{transform}, \texttt{accumulate},
\texttt{all\_of}, \texttt{any\_of}, \texttt{none\_of} i
\texttt{for\_each}. Algorytmy numeryczne znajdują się w nagłówku
\texttt{<numeric>}, pozostałe głównie w \texttt{<algorithm>}.

\begin{lstlisting}[caption={Algorytmy klasyczne i Ranges (C++20); kompletny program.}]
#include <algorithm>
#include <cassert>
#include <numeric>
#include <vector>

int main() {
    std::vector<int> values{1, 2, 3, 4, 5};
    auto even = [](int x) { return x % 2 == 0; };

    assert(std::ranges::count_if(values, even) == 2);
    assert(std::ranges::all_of(values,
        [](int x) { return x > 0; }));
    assert(std::accumulate(values.begin(), values.end(), 0) == 15);

    std::vector<int> squares(values.size());
    std::ranges::transform(values, squares.begin(),
                           [](int x) { return x * x; });
    assert(squares[2] == 9);
}
\end{lstlisting}

\subsection{Sortowanie i stabilność kolejności}

\texttt{std::sort} wymaga iteratorów dostępu swobodnego i ma
złożoność $O(n\log n)$ w najgorszym przypadku (w liczbie porównań).
\texttt{stable\_sort} zachowuje kolejność elementów, które komparator
uznaje za równoważne, zwykle kosztem dodatkowych zasobów; jego
szczegółowe gwarancje złożoności zależą od dostępnej pamięci.
\texttt{nth\_element} umieszcza wybrany element na pozycji, którą
zajmowałby po sortowaniu, bez sortowania całej sekwencji; przeciętna
liczba porównań jest liniowa.

\begin{lstlisting}[caption={Stabilne sortowanie rekordow po kluczu (C++20); kompletny program.}]
#include <algorithm>
#include <cassert>
#include <string>
#include <vector>

struct Record {
    std::string name;
    int group;
};

int main() {
    std::vector<Record> v{{"A", 2}, {"B", 1}, {"C", 2}};
    std::ranges::stable_sort(v, {}, &Record::group);
    assert(v[0].name == "B");
    assert(v[1].name == "A" && v[2].name == "C");
}
\end{lstlisting}

\subsection{Wyszukiwanie binarne i porządek wejścia}

\texttt{binary\_search}, \texttt{lower\_bound}, \texttt{upper\_bound}
i \texttt{equal\_range} wymagają odpowiedniego uporządkowania
(lub partycjonowania względem danego klucza) zakresu. Zastosowanie
do nieuporządkowanych danych daje wynik niespełniający kontraktu
algorytmu. W wersjach na iteratorach nie-losowego dostępu liczba
porównań może pozostać logarytmiczna, lecz przesuwanie iteratorów
może być liniowe.

\begin{lstlisting}[caption={Wyszukiwanie miejsca wstawienia w posortowanym vector (C++11); kompletny program.}]
#include <algorithm>
#include <cassert>
#include <vector>

int main() {
    std::vector<int> sorted{1, 3, 3, 7};
    auto it = std::lower_bound(sorted.begin(), sorted.end(), 3);
    assert(it == sorted.begin() + 1);

    auto upper = std::upper_bound(sorted.begin(), sorted.end(), 3);
    assert(upper == sorted.begin() + 3);
    assert(std::binary_search(sorted.begin(), sorted.end(), 7));
}
\end{lstlisting}

\subsection{Idiom \emph{erase--remove} i \texttt{std::erase\_if}}

\texttt{std::remove} nie usuwa elementów \emph{z kontenera}:
przestawia elementy pozostające i zwraca logiczny nowy koniec.
Rozmiar kontenera zmniejsza dopiero wywołanie \texttt{erase}.
Od C++20 wygodniej używać \texttt{std::erase} i
\texttt{std::erase\_if} dla wspieranych kontenerów.

\begin{lstlisting}[caption={Dwa sposoby usuwania z vector (C++20); kompletny program.}]
#include <algorithm>
#include <cassert>
#include <vector>

int main() {
    std::vector<int> a{1, 2, 3, 4, 5};
    auto last = std::remove_if(a.begin(), a.end(),
                               [](int x) { return x % 2 == 0; });
    a.erase(last, a.end());
    assert((a == std::vector<int>{1, 3, 5}));

    std::vector<int> b{1, 2, 3, 4, 5};
    auto removed = std::erase_if(b,
                                 [](int x) { return x % 2 == 0; });
    assert(removed == 2 && b == a);
}
\end{lstlisting}

Dla kontenerów asocjacyjnych nie wolno stosować klasycznego
\texttt{remove\_if} tak jak dla wektora: ich elementy mają klucze,
których pozycja jest regulowana przez strukturę danych. Stosujemy
\texttt{erase} z iteratorem lub \texttt{std::erase\_if} (C++20).

\subsection{Algorytmy kopca i \texttt{std::priority\_queue}}

\texttt{make\_heap}, \texttt{push\_heap}, \texttt{pop\_heap} i
\texttt{sort\_heap} pozwalają użyć własnego wektora jako kopca.
\texttt{make\_heap} działa w czasie liniowym, a operacje dopisania i
usunięcia maksimum w $O(\log n)$. Są to operacje inne niż pełne sortowanie.

\section{Przypadki projektowe: dobór kontenera do problemu}
\label{sec:cpp13-use-cases}

\subsection{Histogram, deduplikacja i lista wyników}

Załóżmy, że chcemy zebrać wartości, usunąć duplikaty i otrzymać
wynik w kolejności rosnącej. Możemy użyć \texttt{std::set}, albo
\texttt{vector} z \texttt{sort} i \texttt{unique}. Drugi sposób bywa
korzystniejszy w przetwarzaniu wsadowym, ponieważ \texttt{vector}
przechowuje dane ciągle i ogranicza liczbę alokacji węzłów.

\begin{lstlisting}[caption={Deduplikacja przez sort i unique (C++11); kompletny program.}]
#include <algorithm>
#include <cassert>
#include <vector>

int main() {
    std::vector<int> keys{5, 2, 5, 1, 2, 8};
    std::sort(keys.begin(), keys.end());
    keys.erase(std::unique(keys.begin(), keys.end()), keys.end());
    assert((keys == std::vector<int>{1, 2, 5, 8}));
}
\end{lstlisting}

\subsection{Własność elementów w kontenerach polimorficznych}

Przechowywanie obiektów klas pochodnych przez wartość w
\texttt{vector<Base>} może powodować \emph{object slicing}. Dla
polimorfizmu dynamicznego używamy np. wektora wskaźników
\texttt{unique\_ptr<Base>}. Przy zamkniętej rodzinie typów można użyć
wektora z elementami \texttt{variant} (zob. rozdział~\ref{chap:cpp14-type-erasure}).
Różnicę między polimorfizmem a obcięciem obiektu omawia
rozdział~\ref{chap:cpp-object-model}.

\begin{lstlisting}[caption={Kontener wlasciwych obiektow polimorficznych (C++14); kompletny program.}]
#include <cassert>
#include <memory>
#include <vector>

struct Shape {
    virtual ~Shape() = default;
    virtual int area() const = 0;
};
struct Square : Shape {
    int side;
    explicit Square(int x) : side(x) {}
    int area() const override { return side * side; }
};

int main() {
    std::vector<std::unique_ptr<Shape>> shapes;
    shapes.push_back(std::make_unique<Square>(3));
    assert(shapes.front()->area() == 9);
}
\end{lstlisting}

\subsection{Kontenery a współbieżność}

Samo wybranie kontenera standardowego nie czyni jego operacji
modyfikujących bezpiecznymi dla współbieżnego wykonania. Równoczesne
modyfikowanie wspólnego \texttt{vector} przez niezależne
\texttt{push\_back()} bez synchronizacji jest wyścigiem danych.
Można natomiast równolegle odczytywać elementy kontenera, jeśli żaden
wątek ich nie modyfikuje, a czas życia kontenera jest zachowany.
Dla wielu kontenerów można także modyfikować \emph{różne istniejące}
elementy równolegle, jeśli nie dochodzi do konfliktu pamięci; wyjątek
stanowi m.in. specjalizacja \texttt{vector<bool>}. Zmiany strukturalne
wymagają osobnej analizy i synchronizacji (zob. rozdział~\ref{chap:concurrency}).

\section{Złożoność, lokalność pamięci i pomiar wydajności}
\label{sec:cpp13-performance}

Teoretyczna złożoność to jeden z elementów oceny. Istotne są również:
liczba alokacji, liczba odwołań do rozproszonych adresów, rozmiar
elementu, przewidywalność gałęzi i zachowanie cache procesora.
Dla wielu zastosowań \texttt{vector} z wyszukiwaniem liniowym
może być szybszy od struktury drzewiastej przy małej liczbie kluczy.
To nie przeczy złożoności asymptotycznej: stałe i charakter obciążenia
są częścią kosztu rzeczywistego programu.

\subsection{Koszt sumowania ciągłych danych}

\begin{lstlisting}[caption={Prosty szkielet pomiaru vector i list (C++20); kompletny program.}]
#include <chrono>
#include <cstdint>
#include <iostream>
#include <list>
#include <numeric>
#include <utility>
#include <vector>

int main() {
    constexpr int N = 200000;
    std::vector<int> v(N, 1);
    std::list<int> l(v.begin(), v.end());

    auto measure = [](const auto& xs) {
        auto begin = std::chrono::steady_clock::now();
        std::int64_t sum = std::accumulate(xs.begin(), xs.end(),
                                           std::int64_t{0});
        auto end = std::chrono::steady_clock::now();
        auto ns = std::chrono::duration_cast<
            std::chrono::nanoseconds>(end - begin).count();
        return std::pair{sum, ns};
    };

    auto [sum_v, time_v] = measure(v);
    auto [sum_l, time_l] = measure(l);
    std::cout << "vector: " << sum_v << ", " << time_v << " ns\n";
    std::cout << "list:   " << sum_l << ", " << time_l << " ns\n";
}
\end{lstlisting}

To \emph{ilustracja sposobu pomiaru}, a nie miarodajny benchmark.
Jednorazowy wynik bywa zdominowany przez rozgrzanie cache, optymalizacje
kompilatora, scheduler i zakłócenia systemowe. Poprawniejszy eksperyment
powtarza pomiary, rozdziela przygotowanie danych od mierzonej operacji,
uwzględnia pamięć, różne rozmiary danych, medianę i rozrzut pomiarów.
Dla wiarygodnych porównań warto użyć np. Google Benchmark i profilera.

\subsection{Przybliżone reguły wyboru}

\begin{center}
\small
\begin{tabularx}{\textwidth}{@{}lX@{}}
\toprule
\textbf{Wymaganie} & \textbf{Rozsądny wybór początkowy} \\
\midrule
Dużo iteracji, obliczeń, dostęp indeksowany & \texttt{vector}; przy stałym rozmiarze \texttt{array} \\
Częste dodawanie i usuwanie z obu końców & \texttt{deque} \\
Utrzymanie kolejności kluczy, wyszukiwanie przedziałowe & \texttt{map} / \texttt{set} \\
Wyszukiwanie po kluczu bez wymogu kolejności & \texttt{unordered\_map} / \texttt{unordered\_set} \\
Przenoszenie istniejących węzłów i stabilne iteratory & \texttt{list} \\
FIFO / LIFO & \texttt{queue} / \texttt{stack} \\
Wielokrotne pobieranie największego lub najmniejszego & \texttt{priority\_queue} \\
\bottomrule
\end{tabularx}
\end{center}

\section{Najczęstsze błędy i dobre praktyki}
\label{sec:cpp13-pitfalls}

\begin{enumerate}
\item Używanie \texttt{vector::operator[]} bez pewności, że indeks jest
  poprawny; w kodzie granicznym warto stosować \texttt{at()}.
\item Traktowanie \texttt{reserve} jako tworzenia elementów (mylenie
  \texttt{size} z \texttt{capacity}).
\item Korzystanie z iteratorów lub referencji po unieważnieniu przez
  realokację, \texttt{erase} lub \texttt{rehash}.
\item Zakładanie, że \texttt{emplace\_back} zawsze jest szybsze od
  \texttt{push\_back}.
\item Używanie \texttt{map::operator[]} do samego sprawdzania istnienia
  klucza, z niezamierzonym wstawieniem nowej wartości.
\item Pisanie komparatorów naruszających \emph{strict weak ordering};
  komparator \texttt{a <= b} nie jest poprawną relacją porządkującą.
\item Używanie \texttt{std::sort} na \texttt{std::list} lub
  \texttt{std::lower\_bound} na danych bez odpowiedniego porządku.
\item Wnioskowanie o kolejności elementów \texttt{unordered\_map}
  na podstawie przypadkowego przebiegu programu.
\item Niepotrzebne stosowanie \texttt{list} z przekonaniem, że
  wstawienie $O(1)$ zawsze wygrywa z \texttt{vector}.
\item Wprowadzanie równoczesnych modyfikacji kontenera z wielu wątków
  bez uzasadnienia braku konfliktu i odpowiedniej synchronizacji.
\end{enumerate}

\section{Pytania kontrolne}

\begin{enumerate}
\item Jaka jest różnica między \texttt{reserve} a \texttt{resize}, i
  które z nich tworzy nowe elementy?
\item Kiedy \texttt{push\_back} na \texttt{vector} unieważnia wszystkie
  referencje, a kiedy zachowuje referencje do wcześniejszych elementów?
\item Dlaczego \texttt{map::operator[]} może być niepożądany podczas
  sprawdzania istnienia klucza?
\item Jakie wymagania muszą spełniać \texttt{Hash} i \texttt{Eq} w
  \texttt{unordered\_map}?
\item Dlaczego \texttt{std::remove} nie zmienia rozmiaru \texttt{vector}?
\item Jak różnią się koszty i kontrakty \texttt{sort}, \texttt{stable\_sort},
  \texttt{nth\_element} i \texttt{lower\_bound}?
\item W jakich okolicznościach \texttt{list::splice} zapewnia korzyść
  niedostępną dla zwykłego \texttt{vector}?
\end{enumerate}

\section{Podsumowanie rozdziału}

Dobry kod C++ wybiera kontener zgodnie z potrzebnymi operacjami i
własnościami pamięci, a nie jedynie na podstawie pojedynczego rzędu
złożoności. \texttt{vector} to często najlepszy punkt wyjścia;
\texttt{deque}, kontenery drzewiaste, haszujące i listy rozwiązują
bardziej szczegółowe problemy. Praktycznie równie ważne jak złożoności
są reguły unieważniania iteratorów, wymagania komparatorów,
bezpieczeństwo czasu życia i świadome wykorzystywanie algorytmów
standardowych zamiast ręcznie pisanych pętli.

\section*{Literatura i dokumentacja}
\addcontentsline{toc}{section}{Literatura i dokumentacja}
\begin{itemize}
\item \url{https://en.cppreference.com/w/cpp/container} --- kontenery i
  zestawienie reguł unieważniania iteratorów.
\item \url{https://en.cppreference.com/w/cpp/algorithm} --- algorytmy
  biblioteki standardowej i ich kontrakty.
\item \url{https://en.cppreference.com/w/cpp/container/vector} ---
  \texttt{std::vector} i zarządzanie pojemnością.
\item \url{https://en.cppreference.com/w/cpp/container/unordered_map} ---
  struktury haszujące.
\item Scott Meyers, \emph{Effective Modern C++}, O'Reilly, 2014.
\item Nicolai M. Josuttis, \emph{The C++ Standard Library}, wydanie drugie,
  Addison-Wesley, 2012.
\end{itemize}


\chapter{Kontenery STL i algorytmy}
\label{chap:cpp-stl-containers}

\section{Cel i zakres rozdziału}

Kontener przechowuje grupę elementów i pozwala je odnajdywać lub
zmieniać. Najprostsza mapa wyboru wygląda tak: \texttt{vector} jest
listą z dostępem przez indeks, \texttt{map} słownikiem klucz--wartość,
a \texttt{set} zbiorem wartości bez powtórzeń. Jeśli kolejność kluczy
nie ma znaczenia, istnieją także haszujące warianty
\texttt{unordered\_map} i \texttt{unordered\_set}.

Biblioteka standardowa C++ udostępnia gotowe struktury danych, których wybór ma
istotne konsekwencje dla złożoności operacji, lokalności pamięci, bezpieczeństwa
wyjątkowego i poprawności iteratorów. Samo rozróżnienie między
\texttt{std::vector}, \texttt{std::list} i \texttt{std::map} nie wystarcza:
trzeba rozumieć również \emph{własność elementów}, \emph{unieważnianie
iteratorów}, wymagania dotyczące porządku i haszowania oraz relację między
kontenerami a algorytmami. W dalszej części domyślnie używamy C++20.

\paragraph{Główna zasada projektowa.}
Najpierw ustalamy rzeczywiste operacje: czy potrzebujemy dostępu przez
indeks, wyszukiwania po kluczu, zachowania kolejności, stabilnych adresów,
czy częstych wstawek i usunięć. Dopiero potem wybieramy kontener. Prosty
\texttt{std::vector} pozostaje często najlepszym rozwiązaniem, nawet gdy
jego teoretyczna złożoność niektórych operacji jest gorsza, ponieważ dane
leżą w sąsiadujących komórkach pamięci.

\section{Klasyfikacja kontenerów i podstawowe pojęcia}
\label{sec:cpp13-taxonomy}

Biblioteka standardowa rozróżnia m.in. kontenery sekwencyjne
(\texttt{array}, \texttt{vector}, \texttt{deque}, \texttt{list},
\texttt{forward\_list}), uporządkowane asocjacyjne (\texttt{map},
\texttt{set} i ich warianty \texttt{multi}) oraz nieuporządkowane
asocjacyjne (\texttt{unordered\_map}, \texttt{unordered\_set} i warianty
\texttt{multi}). \texttt{stack}, \texttt{queue} i \texttt{priority\_queue}
są \emph{adapterami kontenerów}: oferują inny interfejs dostępu do
wewnętrznej struktury danych. \texttt{std::string} jest wyspecjalizowaną
strukturą sekwencyjną, ale semantykę tekstu omawiamy oddzielnie.

\begin{center}
\small
\begin{tabularx}{\textwidth}{@{}lXXX@{}}
\toprule
\textbf{Kontener} & \textbf{Dostęp} & \textbf{Wstawienie / usunięcie} & \textbf{Typowe zastosowanie} \\
\midrule
\texttt{vector} & Indeks $O(1)$ & Koniec: zamortyzowane $O(1)$; środek: $O(n)$ & Dane zwarte, iteracja, obliczenia \\
\texttt{deque} & Indeks $O(1)$ & Oba końce: $O(1)$; środek: $O(n)$ & Kolejka dwustronna \\
\texttt{list} & Iteracja $O(n)$ & $O(1)$ przy znanym iteratorze & Stabilne węzły, \texttt{splice} \\
\texttt{forward\_list} & Iteracja $O(n)$ & $O(1)$ za poprzednikiem & Lista jednokierunkowa \\
\texttt{map} / \texttt{set} & Szukanie $O(\log n)$ & $O(\log n)$ & Klucze w ustalonej kolejności \\
\texttt{unordered\_map} & Średnio $O(1)$ & Średnio $O(1)$ & Słownik, szybkie wyszukiwanie \\
\bottomrule
\end{tabularx}
\end{center}

W tabeli $n$ oznacza liczbę elementów, a koszty zależą od wymagań
poszczególnych operacji. \texttt{unordered\_map} może w najgorszym
przypadku wykonywać operacje w czasie $O(n)$ (np. przy wielu kolizjach).
Sformułowanie ,,\emph{zamortyzowane} $O(1)$'' oznacza, że pojedyncza
operacja może być kosztowna, ale średni koszt w długiej sekwencji operacji
jest ograniczony stałą.

\paragraph{Wartości, wskaźniki, referencje.}
Kontenery zwykle przechowują \emph{własne obiekty elementów}.
\texttt{std::vector<T>} nie jest kontenerem referencji do obiektów
\texttt{T}; może za to przechowywać np. \texttt{std::reference\_wrapper<T>},
\texttt{T*} lub \texttt{std::unique\_ptr<T>}. Wybór wpływa na własność
obiektów, kopiowanie, czas życia i poprawność dostępu.

\section{\texttt{std::vector}: domyślny kontener sekwencyjny}

\begin{center}\small
\begin{tabularx}{\textwidth}{@{}p{0.34\textwidth}X@{}}
\toprule
Operacja & Za co odpowiada \\
\midrule
\code{size()}, \code{empty()} & Podają liczbę elementów i informację, czy kontener jest pusty. \\
\code{capacity()} & Podaje liczbę elementów, które zmieszczą się bez ponownej alokacji. \\
\code{reserve(n)} & Rezerwuje miejsce na co najmniej \code{n} elementów; nie zmienia \code{size()}. \\
\code{resize(n)} & Zmienia liczbę elementów: tworzy nowe albo niszczy nadmiarowe. \\
\code{push_back(x)} & Dodaje istniejącą wartość na koniec przez kopiowanie lub przeniesienie. \\
\code{emplace_back(args...)} & Konstruuje nowy element na końcu z podanych argumentów. \\
\code{erase(pos)} & Usuwa element; kolejne elementy przesuwają się, a iteratory od miejsca usunięcia tracą ważność. \\
\bottomrule
\end{tabularx}
\end{center}
\label{sec:cpp13-vector}

\texttt{std::vector<T>} przechowuje elementy \emph{ciągle} w pamięci,
co pozwala przekazywać je do interfejsów oczekujących tablicy (np. przez
\texttt{data()} lub \texttt{std::span}). Obiekt wektora zarządza własną
pamięcią i niszczy elementy po zakończeniu swojego czasu życia.

\begin{lstlisting}[caption={Podstawowe operacje na vector (C++11); kompletny program.}]
#include <cassert>
#include <vector>

int main() {
    std::vector<int> v{4, 8, 15};
    v.push_back(16);           // kopiowanie lub przenoszenie argumentu
    v.emplace_back(23);        // konstruowanie nowego elementu

    assert(v.size() == 5);
    assert(v[0] == 4);
    assert(v.at(4) == 23);
    assert(v.front() == 4 && v.back() == 23);

    v.pop_back();
    assert(v.size() == 4);
    v.clear();
    assert(v.empty());
}
\end{lstlisting}

\texttt{operator[]} \emph{nie sprawdza} granic: użycie indeksu spoza zakresu
prowadzi do niezdefiniowanego zachowania. Metoda \texttt{at()} sprawdza
zakres i zgłasza \texttt{std::out\_of\_range}. Metody \texttt{front()},
\texttt{back()} i \texttt{pop\_back()} wymagają niepustego kontenera.

\subsection{\texttt{size}, \texttt{capacity}, \texttt{reserve} i \texttt{resize}}

\texttt{size()} oznacza liczbę \emph{istniejących} elementów, a
\texttt{capacity()} liczbę elementów, dla których przydzielono miejsce
bez konieczności realokacji. \texttt{reserve(k)} zmienia rezerwację
pamięci, nie tworzy elementów; \texttt{resize(k)} zmienia ich liczbę.

\begin{lstlisting}[caption={Rezerwacja pamięci nie zwiększa rozmiaru (C++11); kompletny program.}]
#include <cassert>
#include <vector>

int main() {
    std::vector<int> data;
    data.reserve(100);
    assert(data.size() == 0);
    assert(data.capacity() >= 100);

    data.push_back(7);
    data.resize(4);          // dopisuje trzy elementy o wartosci 0
    assert(data.size() == 4);
    assert(data[0] == 7 && data[3] == 0);
}
\end{lstlisting}

Nie należy wywoływać \texttt{reserve(size()+1)} przed każdym
\texttt{push\_back()}: wymuszenie niemal liniowego wzrostu pojemności
może zniweczyć politykę geometrycznego powiększania bufora i pogorszyć
wydajność. \texttt{shrink\_to\_fit()} jest tylko żądaniem, a nie gwarancją
zmniejszenia pojemności.

\subsection{Realokacja i unieważnianie iteratorów}

Po realokacji \texttt{vector} zmienia adresy wszystkich elementów:
tracą ważność iteratory, referencje i wskaźniki do elementów, również
poprzednia wartość \texttt{data()}. Bez realokacji \texttt{push\_back()}
unieważnia iterator \texttt{end()}, lecz nie unieważnia iteratorów i
referencji do wcześniej istniejących elementów. Wstawienie w środku
bez realokacji unieważnia iteratory i referencje do pozycji wstawienia
oraz wszystkich późniejszych pozycji (włącznie z dawnym \texttt{end()}).

\begin{lstlisting}[caption={Bezpieczne ponowne pobranie iteratora po wstawieniu; kompletny program.}]
#include <cassert>
#include <vector>

int main() {
    std::vector<int> v{1, 3, 4};
    auto pos = v.begin() + 1;
    pos = v.insert(pos, 2);   // przyjmuje zwrocony iterator
    assert(*pos == 2);
    assert((v == std::vector<int>{1, 2, 3, 4}));
}
\end{lstlisting}

\paragraph{Przenoszenie elementów i bezpieczeństwo wyjątkowe.}
Podczas realokacji \texttt{vector} może przenosić lub kopiować elementy.
Dla pewnych typów i operacji warunki silnej gwarancji wyjątkowej zależą
od tego, czy konstruktor przenoszący jest \texttt{noexcept} oraz czy
dostępne jest kopiowanie. Zagadnienie łączy się bezpośrednio z semantyką
przenoszenia omówioną w rozdziale 6.

\subsection{\texttt{push\_back} kontra \texttt{emplace\_back}}

\texttt{emplace\_back(args...)} przekazuje argumenty bezpośrednio do
konstruktora nowego elementu. Nie oznacza to jednak automatycznie
,,zawsze szybciej'': jeśli mamy już obiekt, \texttt{push\_back(std::move(x))}
bywa równie naturalne; przy realokacji nadal konieczne może być
przeniesienie innych elementów. Obie funkcje mogą zgłaszać wyjątki.

\begin{lstlisting}[caption={Konstruowanie elementu z argumentow (C++17); kompletny program.}]
#include <cassert>
#include <string>
#include <utility>
#include <vector>

struct Person {
    std::string name;
    int age;
    Person(std::string n, int a) : name(std::move(n)), age(a) {}
};

int main() {
    std::vector<Person> people;
    people.emplace_back("Ada", 31);
    people.push_back(Person{"Alan", 42});
    assert(people[1].age == 42);
}
\end{lstlisting}

\section{\texttt{std::array}, \texttt{std::deque} i listy}
\label{sec:cpp13-sequential}

\subsection{\texttt{std::array<T,N>}: rozmiar znany w kompilacji}

\texttt{std::array} przechowuje dokładnie $N$ elementów. Nie dokonuje
osobnej dynamicznej alokacji dla swojej zawartości; pamięć elementów
stanowi część samego obiektu \texttt{array}. Można używać iteratorów,
\texttt{size()}, \texttt{data()} i algorytmów, analogicznie do wektora.

\begin{lstlisting}[caption={std::array i algorytmy standardowe (C++17); kompletny program.}]
#include <algorithm>
#include <array>
#include <cassert>
#include <tuple>

int main() {
    std::array<int, 4> x{5, 1, 4, 2};
    std::sort(x.begin(), x.end());
    assert((x == std::array<int, 4>{1, 2, 4, 5}));
    static_assert(std::tuple_size_v<decltype(x)> == 4);
}
\end{lstlisting}

\subsection{\texttt{std::deque}: dwa tanie końce}

\texttt{deque} udostępnia dostęp indeksowany w czasie $O(1)$, lecz
zwykle nie jest jednym ciągłym blokiem pamięci. Wstawianie i usuwanie
na początku i końcu ma stałą złożoność; w przeciwieństwie do
\texttt{vector} ma operacje \texttt{push\_front()} i \texttt{pop\_front()}.
\texttt{deque} nie udostępnia \texttt{data()} dla wszystkich elementów.

\begin{lstlisting}[caption={Kolejka dwustronna (C++11); kompletny program.}]
#include <cassert>
#include <deque>

int main() {
    std::deque<int> q{2, 3};
    q.push_front(1);
    q.push_back(4);
    assert(q[0] == 1 && q[3] == 4);
    q.pop_front();
    assert(q.front() == 2);
}
\end{lstlisting}

Wstawienie na końcu \texttt{deque} może unieważnić \emph{iteratory},
ale zachowuje referencje do istniejących elementów. Szczegółowe reguły
różnią się dla wstawień i usunięć w środku oraz na końcach; nie należy
mechanicznie przenosić reguł \texttt{vector} na \texttt{deque}.

\subsection{\texttt{std::list} i \texttt{std::forward\_list}}

\texttt{list} jest listą dwukierunkową. Wstawianie i usuwanie przy
znanym iteratorze odbywa się w stałym czasie; iteratory i referencje
do pozostałych elementów pozostają ważne. Wyszukanie elementu przez
wartość wymaga jednak przejścia listy $O(n)$, podobnie jak uzyskanie
elementu o zadanym indeksie. W przeciwieństwie do \texttt{vector},
\texttt{list} nie oferuje operatora \texttt{[]} ani iteratorów dostępu
swobodnego. Do sortowania używamy \texttt{list::sort()}, a nie
\texttt{std::sort()}.

\begin{lstlisting}[caption={Przeniesienie wezla przez list::splice (C++11); kompletny program.}]
#include <cassert>
#include <iterator>
#include <list>

int main() {
    std::list<int> source{1, 2, 3};
    std::list<int> target{10};
    auto it = std::next(source.begin()); // wskazuje 2

    target.splice(target.end(), source, it);
    assert(*it == 2); // iterator pozostaje wazny
    assert(source.size() == 2 && target.size() == 2);
    target.sort();
    assert(target.front() == 2);
}
\end{lstlisting}

\texttt{forward\_list} przechowuje węzły listy jednokierunkowej;
wstawianie i usuwanie wykonuje \emph{za} wskazanym poprzednikiem
(\texttt{insert\_after}, \texttt{erase\_after}), a \texttt{before\_begin()}
udostępnia pozycję przed pierwszym elementem. Kontener nie posiada
metody \texttt{size()} (rozmiar można policzyć liniowo przez
\texttt{std::distance}).

\paragraph{Kiedy lista jest rzeczywiście potrzebna?}
Lista bywa przydatna, gdy musimy często \emph{przenosić węzły bez
kopiowania obiektów} i wymagamy stabilnych iteratorów. Dla częstej
iteracji i obliczeń numerycznych \texttt{vector} z reguły zapewnia
lepszą lokalność pamięci. Wybór należy jednak potwierdzić pomiarem.

\section{Kontenery uporządkowane: \texttt{set} i \texttt{map}}
\label{sec:cpp13-ordered}

\texttt{std::set<Key>} przechowuje unikalne klucze, a
\texttt{std::map<Key,T>} unikalne pary klucz--wartość. Elementy są
iterowane w kolejności zdefiniowanej przez komparator; wyszukiwanie,
wstawianie i usuwanie są logarytmiczne. Warianty
\texttt{multiset}/\texttt{multimap} dopuszczają powtarzające się klucze.

\subsection{\texttt{map::operator[]}, \texttt{at} i \texttt{find}}

Odwołanie \texttt{m[key]} do nieistniejącego klucza \emph{wstawia} nowy
element z wartością inicjalizowaną domyślnie; może zatem modyfikować
strukturę danych. \texttt{at()} zgłasza wyjątek, jeżeli klucz nie
istnieje. \texttt{find()} zwraca iterator, a \texttt{contains()} (C++20)
zwraca wartość logiczną bez wstawiania.

\begin{lstlisting}[caption={Liczenie czestosci slow z map (C++20); kompletny program.}]
#include <cassert>
#include <map>
#include <string>
#include <vector>

int main() {
    std::vector<std::string> words{"kot", "pies", "kot"};
    std::map<std::string, int> count;
    for (const auto& word : words) ++count[word];

    assert(count.at("kot") == 2);
    assert(count.contains("pies"));
    assert(!count.contains("lis"));
    assert(count.find("lis") == count.end());
}
\end{lstlisting}

\subsection{\texttt{insert}, \texttt{try\_emplace}, \texttt{insert\_or\_assign}}

\texttt{insert} podejmuje próbę dodania gotowej pary.
\texttt{try\_emplace(key,args...)} (C++17) konstruuje wartość
w miejscu z przekazanych argumentów tylko wtedy, gdy klucz nie występuje.
Ważne zastrzeżenie: \emph{wyrażenia będące argumentami wywołania}
są obliczane jeszcze przed wejściem do \texttt{try\_emplace}; nie
należy więc podawać niepotrzebnie kosztownego wyrażenia tworzącego
tymczasowy obiekt. \texttt{insert\_or\_assign} dodaje lub zastępuje wartość.

\begin{lstlisting}[caption={Rozne operacje modyfikacji map (C++17); kompletny program.}]
#include <cassert>
#include <map>
#include <memory>
#include <string>

struct Heavy {
    static inline int constructions = 0; // C++17
    int value;
    explicit Heavy(int x) : value(x) { ++constructions; }
};

int main() {
    std::map<std::string, Heavy> m;
    auto [it, added] = m.try_emplace("a", 3);
    assert(added && it->second.value == 3);

    auto [again, added_again] = m.try_emplace("a", 4);
    assert(!added_again && again->second.value == 3);
    assert(Heavy::constructions == 1);

    m.insert_or_assign("a", Heavy{9});
    assert(m.at("a").value == 9);

    std::map<int, std::unique_ptr<int>> owned;
    owned.try_emplace(1, std::make_unique<int>(10));
    assert(*owned.at(1) == 10);
}
\end{lstlisting}

\subsection{Własny komparator i ścisły słaby porządek}

Porównanie musi spełniać warunek \emph{strict weak ordering}; w
szczególności komparator powinien być irrefleksyjny
($\neg\operatorname{comp}(x,x)$) i przechodni. Jeżeli dwie wartości
$ a,b $ nie spełniają ani \texttt{comp(a,b)}, ani \texttt{comp(b,a)},
kontener uznaje je za równoważne jako klucze --- nawet gdy
\texttt{operator==} zwraca \texttt{false}. Wymagania tego porządku
są istotne również dla \texttt{sort} i \texttt{lower\_bound}.
Na przykład komparator ignorujący wielkość liter uzna \texttt{"ABC"}
i \texttt{"abc"} za ten sam klucz: druga próba wstawienia do
\texttt{map} nie utworzy nowego elementu.

\begin{lstlisting}[caption={Mapowanie bez rozrozniania wielkosci liter: prosty komparator; kompletny program.}]
#include <cassert>
#include <cctype>
#include <map>
#include <string>

struct CaseInsensitiveLess {
    bool operator()(const std::string& a,
                    const std::string& b) const {
        auto ai = a.begin();
        auto bi = b.begin();
        while (ai != a.end() && bi != b.end()) {
            auto x = static_cast<unsigned char>(*ai++);
            auto y = static_cast<unsigned char>(*bi++);
            int lx = std::tolower(x);
            int ly = std::tolower(y);
            if (lx != ly) return lx < ly;
        }
        return a.size() < b.size();
    }
};

int main() {
    std::map<std::string, int, CaseInsensitiveLess> m;
    m.emplace("ABC", 1);
    m.emplace("abc", 2); // rownowazny klucz
    assert(m.size() == 1 && m.begin()->second == 1);
}
\end{lstlisting}

Przykład porównuje bajty w podstawowym środowisku znakowym;
\texttt{std::tolower} z bieżącą lokalizacją C nie realizuje ogólnego
Unicode case folding. Do pełnego porównywania napisów w Unicode trzeba
zastosować osobne, świadomie wybrane reguły normalizacji i kolacji.

\subsection{Wyszukiwanie heterogeniczne}

Jeżeli komparator mapy jest \emph{transparentny} (np. \texttt{std::less<>}),
możemy szukać w mapie o kluczach \texttt{std::string} za pomocą
\texttt{std::string\_view}, bez tworzenia tymczasowego ciągu znaków.
Heterogeniczne przeciążenia \texttt{find} są dostępne
od C++14, a \texttt{contains} od C++20.

\begin{lstlisting}[caption={Heterogeniczne wyszukiwanie w map (C++20); kompletny program.}]
#include <cassert>
#include <functional>
#include <map>
#include <string>
#include <string_view>

int main() {
    std::map<std::string, int, std::less<>> m{{"alpha", 3}};
    std::string_view key = "alpha";
    auto it = m.find(key);
    assert(it != m.end() && it->second == 3);
    assert(m.contains(key));
}
\end{lstlisting}

\subsection{Zmiana klucza przez uchwyt węzła (C++17)}

Kluczy w \texttt{map} nie wolno modyfikować bezpośrednio przez iterator.
Uchwyt \texttt{node\_handle} uzyskany przez \texttt{extract()} pozwala
wyjąć węzeł, zmienić jego klucz i wstawić z powrotem. Trzeba pamiętać,
że ponowne wstawienie może się nie udać z powodu kolizji kluczy.

\begin{lstlisting}[caption={Zmiana klucza przez extract (C++17); kompletny program.}]
#include <cassert>
#include <map>
#include <string>
#include <utility>

int main() {
    std::map<std::string, int> m{{"old", 7}};
    auto node = m.extract("old");
    assert(!node.empty());
    node.key() = "new";
    auto result = m.insert(std::move(node));
    assert(result.inserted && m.at("new") == 7);
}
\end{lstlisting}

\section{Kontenery haszujące: \texttt{unordered\_map} i \texttt{unordered\_set}}
\label{sec:cpp13-unordered}

Kontener \texttt{unordered\_map<Key,T,Hash,Eq>} rozmieszcza elementy
w kubełkach. Funkcja \texttt{Hash} przekształca klucz w wartość typu
\texttt{std::size\_t}; \texttt{Eq} rozstrzyga równość kluczy.
Wymagany jest warunek zgodności:
\[
\operatorname{Eq}(a,b) \Longrightarrow
\operatorname{Hash}(a)=\operatorname{Hash}(b).
\]
Warunek odwrotny nie zachodzi: kolizje haszy są dozwolone i
rozstrzygane przez porównanie kluczy. Iteracja po kontenerze
nie ma stabilnego porządku i nie należy z niego korzystać jako z
interfejsu logicznego programu.
Można wyobrazić sobie hasz jako wskazówkę, w którym kubełku szukać.
Jeśli dwa różne klucze trafią do tego samego kubełka, \texttt{Eq}
sprawdza, czy są naprawdę równe. Stała w poniższym \texttt{PointHash}
służy jedynie do wymieszania haszy obu współrzędnych.

\begin{lstlisting}[caption={Wlasny hash i porownanie kluczy (C++20); kompletny program.}]
#include <cassert>
#include <cstddef>
#include <functional>
#include <unordered_map>

struct Point {
    int x, y;
    friend bool operator==(const Point&, const Point&) = default;
};

struct PointHash {
    std::size_t operator()(const Point& p) const noexcept {
        std::size_t h1 = std::hash<int>{}(p.x);
        std::size_t h2 = std::hash<int>{}(p.y);
        return h1 ^ (h2 + std::size_t{0x9e3779b9}
                        + (h1 << 6) + (h1 >> 2));
    }
};

int main() {
    std::unordered_map<Point, int, PointHash> m;
    m[{2, 5}] = 42;
    assert(m.at(Point{2, 5}) == 42);
}
\end{lstlisting}

Ta funkcja mieszająca służy demonstracji składni, a nie zastosowaniom
kryptograficznym. Dla publicznie sterowanych kluczy przy projektowaniu
systemów wystawionych na ataki kolizyjne należy rozważyć również
odporność na złośliwe dane wejściowe.

\subsection{\texttt{load\_factor}, \texttt{reserve} i \texttt{rehash}}

\texttt{load\_factor()} to stosunek liczby elementów do liczby kubełków.
\texttt{max\_load\_factor()} podaje lub ustawia pożądany maksymalny
współczynnik. \texttt{reserve(n)} żąda przygotowania na $n$ elementów,
a \texttt{rehash(k)} na przynajmniej $k$ kubełków, z uwzględnieniem
limitów wynikających z współczynnika wypełnienia.

\begin{lstlisting}[caption={Rezerwowanie kubelek przed licznymi wstawkami (C++20); kompletny program.}]
#include <cassert>
#include <unordered_set>

int main() {
    std::unordered_set<int> seen;
    seen.reserve(1000);
    for (int i = 0; i < 1000; ++i) seen.insert(i);
    assert(seen.size() == 1000);
    assert(seen.contains(17));
}
\end{lstlisting}

Rehashing unieważnia iteratory, lecz \emph{nie unieważnia} wskaźników
i referencji do elementów. Wstawienie może unieważnić iteratory, jeśli
wywoła rehash. Usunięcie unieważnia iteratory i referencje do usuniętych
elementów. Nie wolno modyfikować kluczy w sposób zaburzający ich hash
lub równość podczas przebywania elementów w kontenerze.

\section{Adaptery: \texttt{stack}, \texttt{queue}, \texttt{priority\_queue}}
\label{sec:cpp13-adaptors}

\texttt{stack} udostępnia schemat LIFO, \texttt{queue} --- FIFO,
a \texttt{priority\_queue} --- dostęp do elementu o najwyższym
priorytecie według komparatora (domyślnie do największego).
Adaptery zwykle nie zapewniają bezpośredniej iteracji po elementach.

\begin{lstlisting}[caption={Trzy adaptery standardowe (C++11); kompletny program.}]
#include <cassert>
#include <functional>
#include <queue>
#include <stack>
#include <vector>

int main() {
    std::stack<int> s;
    s.push(1); s.push(2);
    assert(s.top() == 2);
    s.pop();

    std::queue<int> q;
    q.push(1); q.push(2);
    assert(q.front() == 1);
    q.pop();

    std::priority_queue<int, std::vector<int>, std::greater<int>> pq;
    pq.push(7); pq.push(2); pq.push(5);
    assert(pq.top() == 2); // kopiec minimum
}
\end{lstlisting}

\texttt{priority\_queue} zwykle opiera się na kopcu: \texttt{top()} ma
koszt $O(1)$, a \texttt{push()} i \texttt{pop()} $O(\log n)$. Nie jest
to kolejka FIFO ani struktura z aktualizacją priorytetu dowolnego elementu
przez gotowy publiczny interfejs.

\section{Iteratory, stabilność adresów i unieważnianie}
\label{sec:cpp13-invalidation}

Iterator nie jest ogólnie tożsamy ze wskaźnikiem. Może przechowywać
stan niezbędny do poruszania się po drzewie lub kubełkach. Po operacji
modyfikującej trzeba wiedzieć, czy iterator, referencja i wskaźnik do
elementu pozostają ważne. Typowe reguły podsumowano w tabeli.

\begin{center}
\footnotesize
\begin{tabularx}{\textwidth}{@{}lXX@{}}
\toprule
\textbf{Kontener} & \textbf{Wstawienie} & \textbf{Usunięcie} \\
\midrule
\texttt{vector} & Przy realokacji: wszystkie iteratory i referencje; bez realokacji: pozycja wstawienia i dalej & Pozycja usunięcia i dalej (także dawny \texttt{end()}) \\
\texttt{deque} & Na końcach: iteratory tracą ważność, referencje do elementów zwykle pozostają; w środku: także referencje & Zależy od pozycji; usunięcie w środku unieważnia wszystkie iteratory i referencje \\
\texttt{list} & Istniejące iteratory i referencje zachowane & Tylko usunięte elementy \\
\texttt{map}, \texttt{set} & Istniejące iteratory i referencje zachowane & Tylko usunięte elementy \\
\texttt{unordered\_*} & Iteratory mogą utracić ważność przy rehash; referencje zachowane & Tylko usunięte elementy \\
\bottomrule
\end{tabularx}
\end{center}

Dla \texttt{deque} szczególnie ostrożnie trzeba traktować
\texttt{end()}, którego ważność zależy od rodzaju usunięcia.
W tabeli podano reguły najważniejsze dla praktyki, a nie kompletną
specyfikację każdej przeciążonej operacji. Po operacjach zmieniających
strukturę bezpieczną praktyką jest korzystanie z iteratora zwróconego
przez metodę (np. \texttt{erase}, \texttt{insert}).

\subsection{Bezpieczne usuwanie podczas iteracji}

\begin{lstlisting}[caption={Usuwanie elementow map podczas iteracji (C++20); kompletny program.}]
#include <cassert>
#include <map>

int main() {
    std::map<int, int> m{{1, 2}, {2, 3}, {3, 4}};
    for (auto it = m.begin(); it != m.end();) {
        if (it->second % 2 == 0) it = m.erase(it);
        else ++it;
    }
    assert(m.size() == 1 && m.contains(2));
}
\end{lstlisting}

\paragraph{Szczególny przypadek \texttt{vector<bool>}.}
\texttt{std::vector<bool>} jest specjalizacją oszczędzającą pamięć,
która może reprezentować wartości w postaci bitów. Jej \texttt{operator[]}
nie musi zwracać \texttt{bool\&}, tylko specjalny obiekt proxy. Nie
należy oczekiwać normalnej tablicy obiektów \texttt{bool}, ciągłego
przechowywania pojedynczych \texttt{bool} ani bezpiecznej modyfikacji
odrębnych bitów z różnych wątków bez synchronizacji.

\section{Algorytmy standardowe i wymagania wobec iteratorów}

\begin{center}\small
\begin{tabularx}{\textwidth}{@{}p{0.34\textwidth}X@{}}
\toprule
Algorytm & Działanie i warunki \\
\midrule
\code{std::find} & Szuka pierwszej wartości równej wzorcowi; przy braku zwraca iterator końcowy. \\
\code{std::transform} & Stosuje funkcję do zakresu i zapisuje wyniki przez iterator wyjściowy. \\
\code{std::sort} & Sortuje zakres z iteratorami dostępu swobodnego; nie zachowuje kolejności elementów równoważnych. \\
\code{std::stable_sort} & Sortuje, zachowując wzajemną kolejność elementów równoważnych. \\
\code{std::lower_bound} & Szuka pierwszego elementu nie mniejszego od wzorca; wymaga odpowiednio uporządkowanego zakresu. \\
\code{std::erase_if} & Usuwa z kontenera elementy spełniające predykat; dostępne od C++20 dla wspieranych kontenerów. \\
\bottomrule
\end{tabularx}
\end{center}
\label{sec:cpp13-algorithms}

Algorytmy oddzielają \emph{operację} od \emph{reprezentacji danych}.
Wersje klasyczne przyjmują iteratory, a w C++20 wiele algorytmów
w \texttt{std::ranges} można wywołać bezpośrednio na zakresie. Szczegóły
kategorii iteratorów i widoków omawiamy w rozdziale~\ref{chap:ranges-views-cpp}; tutaj skupiamy
się na doborze algorytmu i jego kontrakcie.

\subsection{Wyszukiwanie, transformacje i agregacje}

Do częstych operacji należą \texttt{find}, \texttt{find\_if},
\texttt{count\_if}, \texttt{transform}, \texttt{accumulate},
\texttt{all\_of}, \texttt{any\_of}, \texttt{none\_of} i
\texttt{for\_each}. Algorytmy numeryczne znajdują się w nagłówku
\texttt{<numeric>}, pozostałe głównie w \texttt{<algorithm>}.

\begin{lstlisting}[caption={Algorytmy klasyczne i Ranges (C++20); kompletny program.}]
#include <algorithm>
#include <cassert>
#include <numeric>
#include <vector>

int main() {
    std::vector<int> values{1, 2, 3, 4, 5};
    auto even = [](int x) { return x % 2 == 0; };

    assert(std::ranges::count_if(values, even) == 2);
    assert(std::ranges::all_of(values,
        [](int x) { return x > 0; }));
    assert(std::accumulate(values.begin(), values.end(), 0) == 15);

    std::vector<int> squares(values.size());
    std::ranges::transform(values, squares.begin(),
                           [](int x) { return x * x; });
    assert(squares[2] == 9);
}
\end{lstlisting}

\subsection{Sortowanie i stabilność kolejności}

\texttt{std::sort} wymaga iteratorów dostępu swobodnego i ma
złożoność $O(n\log n)$ w najgorszym przypadku (w liczbie porównań).
\texttt{stable\_sort} zachowuje kolejność elementów, które komparator
uznaje za równoważne, zwykle kosztem dodatkowych zasobów; jego
szczegółowe gwarancje złożoności zależą od dostępnej pamięci.
\texttt{nth\_element} umieszcza wybrany element na pozycji, którą
zajmowałby po sortowaniu, bez sortowania całej sekwencji; przeciętna
liczba porównań jest liniowa.

\begin{lstlisting}[caption={Stabilne sortowanie rekordow po kluczu (C++20); kompletny program.}]
#include <algorithm>
#include <cassert>
#include <string>
#include <vector>

struct Record {
    std::string name;
    int group;
};

int main() {
    std::vector<Record> v{{"A", 2}, {"B", 1}, {"C", 2}};
    std::ranges::stable_sort(v, {}, &Record::group);
    assert(v[0].name == "B");
    assert(v[1].name == "A" && v[2].name == "C");
}
\end{lstlisting}

\subsection{Wyszukiwanie binarne i porządek wejścia}

\texttt{binary\_search}, \texttt{lower\_bound}, \texttt{upper\_bound}
i \texttt{equal\_range} wymagają odpowiedniego uporządkowania
(lub partycjonowania względem danego klucza) zakresu. Zastosowanie
do nieuporządkowanych danych daje wynik niespełniający kontraktu
algorytmu. W wersjach na iteratorach nie-losowego dostępu liczba
porównań może pozostać logarytmiczna, lecz przesuwanie iteratorów
może być liniowe.

\begin{lstlisting}[caption={Wyszukiwanie miejsca wstawienia w posortowanym vector (C++11); kompletny program.}]
#include <algorithm>
#include <cassert>
#include <vector>

int main() {
    std::vector<int> sorted{1, 3, 3, 7};
    auto it = std::lower_bound(sorted.begin(), sorted.end(), 3);
    assert(it == sorted.begin() + 1);

    auto upper = std::upper_bound(sorted.begin(), sorted.end(), 3);
    assert(upper == sorted.begin() + 3);
    assert(std::binary_search(sorted.begin(), sorted.end(), 7));
}
\end{lstlisting}

\subsection{Idiom \emph{erase--remove} i \texttt{std::erase\_if}}

\texttt{std::remove} nie usuwa elementów \emph{z kontenera}:
przestawia elementy pozostające i zwraca logiczny nowy koniec.
Rozmiar kontenera zmniejsza dopiero wywołanie \texttt{erase}.
Od C++20 wygodniej używać \texttt{std::erase} i
\texttt{std::erase\_if} dla wspieranych kontenerów.

\begin{lstlisting}[caption={Dwa sposoby usuwania z vector (C++20); kompletny program.}]
#include <algorithm>
#include <cassert>
#include <vector>

int main() {
    std::vector<int> a{1, 2, 3, 4, 5};
    auto last = std::remove_if(a.begin(), a.end(),
                               [](int x) { return x % 2 == 0; });
    a.erase(last, a.end());
    assert((a == std::vector<int>{1, 3, 5}));

    std::vector<int> b{1, 2, 3, 4, 5};
    auto removed = std::erase_if(b,
                                 [](int x) { return x % 2 == 0; });
    assert(removed == 2 && b == a);
}
\end{lstlisting}

Dla kontenerów asocjacyjnych nie wolno stosować klasycznego
\texttt{remove\_if} tak jak dla wektora: ich elementy mają klucze,
których pozycja jest regulowana przez strukturę danych. Stosujemy
\texttt{erase} z iteratorem lub \texttt{std::erase\_if} (C++20).

\subsection{Algorytmy kopca i \texttt{std::priority\_queue}}

\texttt{make\_heap}, \texttt{push\_heap}, \texttt{pop\_heap} i
\texttt{sort\_heap} pozwalają użyć własnego wektora jako kopca.
\texttt{make\_heap} działa w czasie liniowym, a operacje dopisania i
usunięcia maksimum w $O(\log n)$. Są to operacje inne niż pełne sortowanie.

\section{Przypadki projektowe: dobór kontenera do problemu}
\label{sec:cpp13-use-cases}

\subsection{Histogram, deduplikacja i lista wyników}

Załóżmy, że chcemy zebrać wartości, usunąć duplikaty i otrzymać
wynik w kolejności rosnącej. Możemy użyć \texttt{std::set}, albo
\texttt{vector} z \texttt{sort} i \texttt{unique}. Drugi sposób bywa
korzystniejszy w przetwarzaniu wsadowym, ponieważ \texttt{vector}
przechowuje dane ciągle i ogranicza liczbę alokacji węzłów.

\begin{lstlisting}[caption={Deduplikacja przez sort i unique (C++11); kompletny program.}]
#include <algorithm>
#include <cassert>
#include <vector>

int main() {
    std::vector<int> keys{5, 2, 5, 1, 2, 8};
    std::sort(keys.begin(), keys.end());
    keys.erase(std::unique(keys.begin(), keys.end()), keys.end());
    assert((keys == std::vector<int>{1, 2, 5, 8}));
}
\end{lstlisting}

\subsection{Własność elementów w kontenerach polimorficznych}

Przechowywanie obiektów klas pochodnych przez wartość w
\texttt{vector<Base>} może powodować \emph{object slicing}. Dla
polimorfizmu dynamicznego używamy np. wektora wskaźników
\texttt{unique\_ptr<Base>}. Przy zamkniętej rodzinie typów można użyć
wektora z elementami \texttt{variant} (zob. rozdział~\ref{chap:cpp14-type-erasure}).
Różnicę między polimorfizmem a obcięciem obiektu omawia
rozdział~\ref{chap:cpp-object-model}.

\begin{lstlisting}[caption={Kontener wlasciwych obiektow polimorficznych (C++14); kompletny program.}]
#include <cassert>
#include <memory>
#include <vector>

struct Shape {
    virtual ~Shape() = default;
    virtual int area() const = 0;
};
struct Square : Shape {
    int side;
    explicit Square(int x) : side(x) {}
    int area() const override { return side * side; }
};

int main() {
    std::vector<std::unique_ptr<Shape>> shapes;
    shapes.push_back(std::make_unique<Square>(3));
    assert(shapes.front()->area() == 9);
}
\end{lstlisting}

\subsection{Kontenery a współbieżność}

Samo wybranie kontenera standardowego nie czyni jego operacji
modyfikujących bezpiecznymi dla współbieżnego wykonania. Równoczesne
modyfikowanie wspólnego \texttt{vector} przez niezależne
\texttt{push\_back()} bez synchronizacji jest wyścigiem danych.
Można natomiast równolegle odczytywać elementy kontenera, jeśli żaden
wątek ich nie modyfikuje, a czas życia kontenera jest zachowany.
Dla wielu kontenerów można także modyfikować \emph{różne istniejące}
elementy równolegle, jeśli nie dochodzi do konfliktu pamięci; wyjątek
stanowi m.in. specjalizacja \texttt{vector<bool>}. Zmiany strukturalne
wymagają osobnej analizy i synchronizacji (zob. rozdział~\ref{chap:concurrency}).

\section{Złożoność, lokalność pamięci i pomiar wydajności}
\label{sec:cpp13-performance}

Teoretyczna złożoność to jeden z elementów oceny. Istotne są również:
liczba alokacji, liczba odwołań do rozproszonych adresów, rozmiar
elementu, przewidywalność gałęzi i zachowanie cache procesora.
Dla wielu zastosowań \texttt{vector} z wyszukiwaniem liniowym
może być szybszy od struktury drzewiastej przy małej liczbie kluczy.
To nie przeczy złożoności asymptotycznej: stałe i charakter obciążenia
są częścią kosztu rzeczywistego programu.

\subsection{Koszt sumowania ciągłych danych}

\begin{lstlisting}[caption={Prosty szkielet pomiaru vector i list (C++20); kompletny program.}]
#include <chrono>
#include <cstdint>
#include <iostream>
#include <list>
#include <numeric>
#include <utility>
#include <vector>

int main() {
    constexpr int N = 200000;
    std::vector<int> v(N, 1);
    std::list<int> l(v.begin(), v.end());

    auto measure = [](const auto& xs) {
        auto begin = std::chrono::steady_clock::now();
        std::int64_t sum = std::accumulate(xs.begin(), xs.end(),
                                           std::int64_t{0});
        auto end = std::chrono::steady_clock::now();
        auto ns = std::chrono::duration_cast<
            std::chrono::nanoseconds>(end - begin).count();
        return std::pair{sum, ns};
    };

    auto [sum_v, time_v] = measure(v);
    auto [sum_l, time_l] = measure(l);
    std::cout << "vector: " << sum_v << ", " << time_v << " ns\n";
    std::cout << "list:   " << sum_l << ", " << time_l << " ns\n";
}
\end{lstlisting}

To \emph{ilustracja sposobu pomiaru}, a nie miarodajny benchmark.
Jednorazowy wynik bywa zdominowany przez rozgrzanie cache, optymalizacje
kompilatora, scheduler i zakłócenia systemowe. Poprawniejszy eksperyment
powtarza pomiary, rozdziela przygotowanie danych od mierzonej operacji,
uwzględnia pamięć, różne rozmiary danych, medianę i rozrzut pomiarów.
Dla wiarygodnych porównań warto użyć np. Google Benchmark i profilera.

\subsection{Przybliżone reguły wyboru}

\begin{center}
\small
\begin{tabularx}{\textwidth}{@{}lX@{}}
\toprule
\textbf{Wymaganie} & \textbf{Rozsądny wybór początkowy} \\
\midrule
Dużo iteracji, obliczeń, dostęp indeksowany & \texttt{vector}; przy stałym rozmiarze \texttt{array} \\
Częste dodawanie i usuwanie z obu końców & \texttt{deque} \\
Utrzymanie kolejności kluczy, wyszukiwanie przedziałowe & \texttt{map} / \texttt{set} \\
Wyszukiwanie po kluczu bez wymogu kolejności & \texttt{unordered\_map} / \texttt{unordered\_set} \\
Przenoszenie istniejących węzłów i stabilne iteratory & \texttt{list} \\
FIFO / LIFO & \texttt{queue} / \texttt{stack} \\
Wielokrotne pobieranie największego lub najmniejszego & \texttt{priority\_queue} \\
\bottomrule
\end{tabularx}
\end{center}

\section{Najczęstsze błędy i dobre praktyki}
\label{sec:cpp13-pitfalls}

\begin{enumerate}
\item Używanie \texttt{vector::operator[]} bez pewności, że indeks jest
  poprawny; w kodzie granicznym warto stosować \texttt{at()}.
\item Traktowanie \texttt{reserve} jako tworzenia elementów (mylenie
  \texttt{size} z \texttt{capacity}).
\item Korzystanie z iteratorów lub referencji po unieważnieniu przez
  realokację, \texttt{erase} lub \texttt{rehash}.
\item Zakładanie, że \texttt{emplace\_back} zawsze jest szybsze od
  \texttt{push\_back}.
\item Używanie \texttt{map::operator[]} do samego sprawdzania istnienia
  klucza, z niezamierzonym wstawieniem nowej wartości.
\item Pisanie komparatorów naruszających \emph{strict weak ordering};
  komparator \texttt{a <= b} nie jest poprawną relacją porządkującą.
\item Używanie \texttt{std::sort} na \texttt{std::list} lub
  \texttt{std::lower\_bound} na danych bez odpowiedniego porządku.
\item Wnioskowanie o kolejności elementów \texttt{unordered\_map}
  na podstawie przypadkowego przebiegu programu.
\item Niepotrzebne stosowanie \texttt{list} z przekonaniem, że
  wstawienie $O(1)$ zawsze wygrywa z \texttt{vector}.
\item Wprowadzanie równoczesnych modyfikacji kontenera z wielu wątków
  bez uzasadnienia braku konfliktu i odpowiedniej synchronizacji.
\end{enumerate}

\section{Pytania kontrolne}

\begin{enumerate}
\item Jaka jest różnica między \texttt{reserve} a \texttt{resize}, i
  które z nich tworzy nowe elementy?
\item Kiedy \texttt{push\_back} na \texttt{vector} unieważnia wszystkie
  referencje, a kiedy zachowuje referencje do wcześniejszych elementów?
\item Dlaczego \texttt{map::operator[]} może być niepożądany podczas
  sprawdzania istnienia klucza?
\item Jakie wymagania muszą spełniać \texttt{Hash} i \texttt{Eq} w
  \texttt{unordered\_map}?
\item Dlaczego \texttt{std::remove} nie zmienia rozmiaru \texttt{vector}?
\item Jak różnią się koszty i kontrakty \texttt{sort}, \texttt{stable\_sort},
  \texttt{nth\_element} i \texttt{lower\_bound}?
\item W jakich okolicznościach \texttt{list::splice} zapewnia korzyść
  niedostępną dla zwykłego \texttt{vector}?
\end{enumerate}

\section{Podsumowanie rozdziału}

Dobry kod C++ wybiera kontener zgodnie z potrzebnymi operacjami i
własnościami pamięci, a nie jedynie na podstawie pojedynczego rzędu
złożoności. \texttt{vector} to często najlepszy punkt wyjścia;
\texttt{deque}, kontenery drzewiaste, haszujące i listy rozwiązują
bardziej szczegółowe problemy. Praktycznie równie ważne jak złożoności
są reguły unieważniania iteratorów, wymagania komparatorów,
bezpieczeństwo czasu życia i świadome wykorzystywanie algorytmów
standardowych zamiast ręcznie pisanych pętli.

\section*{Literatura i dokumentacja}
\addcontentsline{toc}{section}{Literatura i dokumentacja}
\begin{itemize}
\item \url{https://en.cppreference.com/w/cpp/container} --- kontenery i
  zestawienie reguł unieważniania iteratorów.
\item \url{https://en.cppreference.com/w/cpp/algorithm} --- algorytmy
  biblioteki standardowej i ich kontrakty.
\item \url{https://en.cppreference.com/w/cpp/container/vector} ---
  \texttt{std::vector} i zarządzanie pojemnością.
\item \url{https://en.cppreference.com/w/cpp/container/unordered_map} ---
  struktury haszujące.
\item Scott Meyers, \emph{Effective Modern C++}, O'Reilly, 2014.
\item Nicolai M. Josuttis, \emph{The C++ Standard Library}, wydanie drugie,
  Addison-Wesley, 2012.
\end{itemize}


\chapter{Kontenery STL i algorytmy}
\label{chap:cpp-stl-containers}

\section{Cel i zakres rozdziału}

Kontener przechowuje grupę elementów i pozwala je odnajdywać lub
zmieniać. Najprostsza mapa wyboru wygląda tak: \texttt{vector} jest
listą z dostępem przez indeks, \texttt{map} słownikiem klucz--wartość,
a \texttt{set} zbiorem wartości bez powtórzeń. Jeśli kolejność kluczy
nie ma znaczenia, istnieją także haszujące warianty
\texttt{unordered\_map} i \texttt{unordered\_set}.

Biblioteka standardowa C++ udostępnia gotowe struktury danych, których wybór ma
istotne konsekwencje dla złożoności operacji, lokalności pamięci, bezpieczeństwa
wyjątkowego i poprawności iteratorów. Samo rozróżnienie między
\texttt{std::vector}, \texttt{std::list} i \texttt{std::map} nie wystarcza:
trzeba rozumieć również \emph{własność elementów}, \emph{unieważnianie
iteratorów}, wymagania dotyczące porządku i haszowania oraz relację między
kontenerami a algorytmami. W dalszej części domyślnie używamy C++20.

\paragraph{Główna zasada projektowa.}
Najpierw ustalamy rzeczywiste operacje: czy potrzebujemy dostępu przez
indeks, wyszukiwania po kluczu, zachowania kolejności, stabilnych adresów,
czy częstych wstawek i usunięć. Dopiero potem wybieramy kontener. Prosty
\texttt{std::vector} pozostaje często najlepszym rozwiązaniem, nawet gdy
jego teoretyczna złożoność niektórych operacji jest gorsza, ponieważ dane
leżą w sąsiadujących komórkach pamięci.

\section{Klasyfikacja kontenerów i podstawowe pojęcia}
\label{sec:cpp13-taxonomy}

Biblioteka standardowa rozróżnia m.in. kontenery sekwencyjne
(\texttt{array}, \texttt{vector}, \texttt{deque}, \texttt{list},
\texttt{forward\_list}), uporządkowane asocjacyjne (\texttt{map},
\texttt{set} i ich warianty \texttt{multi}) oraz nieuporządkowane
asocjacyjne (\texttt{unordered\_map}, \texttt{unordered\_set} i warianty
\texttt{multi}). \texttt{stack}, \texttt{queue} i \texttt{priority\_queue}
są \emph{adapterami kontenerów}: oferują inny interfejs dostępu do
wewnętrznej struktury danych. \texttt{std::string} jest wyspecjalizowaną
strukturą sekwencyjną, ale semantykę tekstu omawiamy oddzielnie.

\begin{center}
\small
\begin{tabularx}{\textwidth}{@{}lXXX@{}}
\toprule
\textbf{Kontener} & \textbf{Dostęp} & \textbf{Wstawienie / usunięcie} & \textbf{Typowe zastosowanie} \\
\midrule
\texttt{vector} & Indeks $O(1)$ & Koniec: zamortyzowane $O(1)$; środek: $O(n)$ & Dane zwarte, iteracja, obliczenia \\
\texttt{deque} & Indeks $O(1)$ & Oba końce: $O(1)$; środek: $O(n)$ & Kolejka dwustronna \\
\texttt{list} & Iteracja $O(n)$ & $O(1)$ przy znanym iteratorze & Stabilne węzły, \texttt{splice} \\
\texttt{forward\_list} & Iteracja $O(n)$ & $O(1)$ za poprzednikiem & Lista jednokierunkowa \\
\texttt{map} / \texttt{set} & Szukanie $O(\log n)$ & $O(\log n)$ & Klucze w ustalonej kolejności \\
\texttt{unordered\_map} & Średnio $O(1)$ & Średnio $O(1)$ & Słownik, szybkie wyszukiwanie \\
\bottomrule
\end{tabularx}
\end{center}

W tabeli $n$ oznacza liczbę elementów, a koszty zależą od wymagań
poszczególnych operacji. \texttt{unordered\_map} może w najgorszym
przypadku wykonywać operacje w czasie $O(n)$ (np. przy wielu kolizjach).
Sformułowanie ,,\emph{zamortyzowane} $O(1)$'' oznacza, że pojedyncza
operacja może być kosztowna, ale średni koszt w długiej sekwencji operacji
jest ograniczony stałą.

\paragraph{Wartości, wskaźniki, referencje.}
Kontenery zwykle przechowują \emph{własne obiekty elementów}.
\texttt{std::vector<T>} nie jest kontenerem referencji do obiektów
\texttt{T}; może za to przechowywać np. \texttt{std::reference\_wrapper<T>},
\texttt{T*} lub \texttt{std::unique\_ptr<T>}. Wybór wpływa na własność
obiektów, kopiowanie, czas życia i poprawność dostępu.

\section{\texttt{std::vector}: domyślny kontener sekwencyjny}

\begin{center}\small
\begin{tabularx}{\textwidth}{@{}p{0.34\textwidth}X@{}}
\toprule
Operacja & Za co odpowiada \\
\midrule
\code{size()}, \code{empty()} & Podają liczbę elementów i informację, czy kontener jest pusty. \\
\code{capacity()} & Podaje liczbę elementów, które zmieszczą się bez ponownej alokacji. \\
\code{reserve(n)} & Rezerwuje miejsce na co najmniej \code{n} elementów; nie zmienia \code{size()}. \\
\code{resize(n)} & Zmienia liczbę elementów: tworzy nowe albo niszczy nadmiarowe. \\
\code{push_back(x)} & Dodaje istniejącą wartość na koniec przez kopiowanie lub przeniesienie. \\
\code{emplace_back(args...)} & Konstruuje nowy element na końcu z podanych argumentów. \\
\code{erase(pos)} & Usuwa element; kolejne elementy przesuwają się, a iteratory od miejsca usunięcia tracą ważność. \\
\bottomrule
\end{tabularx}
\end{center}
\label{sec:cpp13-vector}

\texttt{std::vector<T>} przechowuje elementy \emph{ciągle} w pamięci,
co pozwala przekazywać je do interfejsów oczekujących tablicy (np. przez
\texttt{data()} lub \texttt{std::span}). Obiekt wektora zarządza własną
pamięcią i niszczy elementy po zakończeniu swojego czasu życia.

\begin{lstlisting}[caption={Podstawowe operacje na vector (C++11); kompletny program.}]
#include <cassert>
#include <vector>

int main() {
    std::vector<int> v{4, 8, 15};
    v.push_back(16);           // kopiowanie lub przenoszenie argumentu
    v.emplace_back(23);        // konstruowanie nowego elementu

    assert(v.size() == 5);
    assert(v[0] == 4);
    assert(v.at(4) == 23);
    assert(v.front() == 4 && v.back() == 23);

    v.pop_back();
    assert(v.size() == 4);
    v.clear();
    assert(v.empty());
}
\end{lstlisting}

\texttt{operator[]} \emph{nie sprawdza} granic: użycie indeksu spoza zakresu
prowadzi do niezdefiniowanego zachowania. Metoda \texttt{at()} sprawdza
zakres i zgłasza \texttt{std::out\_of\_range}. Metody \texttt{front()},
\texttt{back()} i \texttt{pop\_back()} wymagają niepustego kontenera.

\subsection{\texttt{size}, \texttt{capacity}, \texttt{reserve} i \texttt{resize}}

\texttt{size()} oznacza liczbę \emph{istniejących} elementów, a
\texttt{capacity()} liczbę elementów, dla których przydzielono miejsce
bez konieczności realokacji. \texttt{reserve(k)} zmienia rezerwację
pamięci, nie tworzy elementów; \texttt{resize(k)} zmienia ich liczbę.

\begin{lstlisting}[caption={Rezerwacja pamięci nie zwiększa rozmiaru (C++11); kompletny program.}]
#include <cassert>
#include <vector>

int main() {
    std::vector<int> data;
    data.reserve(100);
    assert(data.size() == 0);
    assert(data.capacity() >= 100);

    data.push_back(7);
    data.resize(4);          // dopisuje trzy elementy o wartosci 0
    assert(data.size() == 4);
    assert(data[0] == 7 && data[3] == 0);
}
\end{lstlisting}

Nie należy wywoływać \texttt{reserve(size()+1)} przed każdym
\texttt{push\_back()}: wymuszenie niemal liniowego wzrostu pojemności
może zniweczyć politykę geometrycznego powiększania bufora i pogorszyć
wydajność. \texttt{shrink\_to\_fit()} jest tylko żądaniem, a nie gwarancją
zmniejszenia pojemności.

\subsection{Realokacja i unieważnianie iteratorów}

Po realokacji \texttt{vector} zmienia adresy wszystkich elementów:
tracą ważność iteratory, referencje i wskaźniki do elementów, również
poprzednia wartość \texttt{data()}. Bez realokacji \texttt{push\_back()}
unieważnia iterator \texttt{end()}, lecz nie unieważnia iteratorów i
referencji do wcześniej istniejących elementów. Wstawienie w środku
bez realokacji unieważnia iteratory i referencje do pozycji wstawienia
oraz wszystkich późniejszych pozycji (włącznie z dawnym \texttt{end()}).

\begin{lstlisting}[caption={Bezpieczne ponowne pobranie iteratora po wstawieniu; kompletny program.}]
#include <cassert>
#include <vector>

int main() {
    std::vector<int> v{1, 3, 4};
    auto pos = v.begin() + 1;
    pos = v.insert(pos, 2);   // przyjmuje zwrocony iterator
    assert(*pos == 2);
    assert((v == std::vector<int>{1, 2, 3, 4}));
}
\end{lstlisting}

\paragraph{Przenoszenie elementów i bezpieczeństwo wyjątkowe.}
Podczas realokacji \texttt{vector} może przenosić lub kopiować elementy.
Dla pewnych typów i operacji warunki silnej gwarancji wyjątkowej zależą
od tego, czy konstruktor przenoszący jest \texttt{noexcept} oraz czy
dostępne jest kopiowanie. Zagadnienie łączy się bezpośrednio z semantyką
przenoszenia omówioną w rozdziale 6.

\subsection{\texttt{push\_back} kontra \texttt{emplace\_back}}

\texttt{emplace\_back(args...)} przekazuje argumenty bezpośrednio do
konstruktora nowego elementu. Nie oznacza to jednak automatycznie
,,zawsze szybciej'': jeśli mamy już obiekt, \texttt{push\_back(std::move(x))}
bywa równie naturalne; przy realokacji nadal konieczne może być
przeniesienie innych elementów. Obie funkcje mogą zgłaszać wyjątki.

\begin{lstlisting}[caption={Konstruowanie elementu z argumentow (C++17); kompletny program.}]
#include <cassert>
#include <string>
#include <utility>
#include <vector>

struct Person {
    std::string name;
    int age;
    Person(std::string n, int a) : name(std::move(n)), age(a) {}
};

int main() {
    std::vector<Person> people;
    people.emplace_back("Ada", 31);
    people.push_back(Person{"Alan", 42});
    assert(people[1].age == 42);
}
\end{lstlisting}

\section{\texttt{std::array}, \texttt{std::deque} i listy}
\label{sec:cpp13-sequential}

\subsection{\texttt{std::array<T,N>}: rozmiar znany w kompilacji}

\texttt{std::array} przechowuje dokładnie $N$ elementów. Nie dokonuje
osobnej dynamicznej alokacji dla swojej zawartości; pamięć elementów
stanowi część samego obiektu \texttt{array}. Można używać iteratorów,
\texttt{size()}, \texttt{data()} i algorytmów, analogicznie do wektora.

\begin{lstlisting}[caption={std::array i algorytmy standardowe (C++17); kompletny program.}]
#include <algorithm>
#include <array>
#include <cassert>
#include <tuple>

int main() {
    std::array<int, 4> x{5, 1, 4, 2};
    std::sort(x.begin(), x.end());
    assert((x == std::array<int, 4>{1, 2, 4, 5}));
    static_assert(std::tuple_size_v<decltype(x)> == 4);
}
\end{lstlisting}

\subsection{\texttt{std::deque}: dwa tanie końce}

\texttt{deque} udostępnia dostęp indeksowany w czasie $O(1)$, lecz
zwykle nie jest jednym ciągłym blokiem pamięci. Wstawianie i usuwanie
na początku i końcu ma stałą złożoność; w przeciwieństwie do
\texttt{vector} ma operacje \texttt{push\_front()} i \texttt{pop\_front()}.
\texttt{deque} nie udostępnia \texttt{data()} dla wszystkich elementów.

\begin{lstlisting}[caption={Kolejka dwustronna (C++11); kompletny program.}]
#include <cassert>
#include <deque>

int main() {
    std::deque<int> q{2, 3};
    q.push_front(1);
    q.push_back(4);
    assert(q[0] == 1 && q[3] == 4);
    q.pop_front();
    assert(q.front() == 2);
}
\end{lstlisting}

Wstawienie na końcu \texttt{deque} może unieważnić \emph{iteratory},
ale zachowuje referencje do istniejących elementów. Szczegółowe reguły
różnią się dla wstawień i usunięć w środku oraz na końcach; nie należy
mechanicznie przenosić reguł \texttt{vector} na \texttt{deque}.

\subsection{\texttt{std::list} i \texttt{std::forward\_list}}

\texttt{list} jest listą dwukierunkową. Wstawianie i usuwanie przy
znanym iteratorze odbywa się w stałym czasie; iteratory i referencje
do pozostałych elementów pozostają ważne. Wyszukanie elementu przez
wartość wymaga jednak przejścia listy $O(n)$, podobnie jak uzyskanie
elementu o zadanym indeksie. W przeciwieństwie do \texttt{vector},
\texttt{list} nie oferuje operatora \texttt{[]} ani iteratorów dostępu
swobodnego. Do sortowania używamy \texttt{list::sort()}, a nie
\texttt{std::sort()}.

\begin{lstlisting}[caption={Przeniesienie wezla przez list::splice (C++11); kompletny program.}]
#include <cassert>
#include <iterator>
#include <list>

int main() {
    std::list<int> source{1, 2, 3};
    std::list<int> target{10};
    auto it = std::next(source.begin()); // wskazuje 2

    target.splice(target.end(), source, it);
    assert(*it == 2); // iterator pozostaje wazny
    assert(source.size() == 2 && target.size() == 2);
    target.sort();
    assert(target.front() == 2);
}
\end{lstlisting}

\texttt{forward\_list} przechowuje węzły listy jednokierunkowej;
wstawianie i usuwanie wykonuje \emph{za} wskazanym poprzednikiem
(\texttt{insert\_after}, \texttt{erase\_after}), a \texttt{before\_begin()}
udostępnia pozycję przed pierwszym elementem. Kontener nie posiada
metody \texttt{size()} (rozmiar można policzyć liniowo przez
\texttt{std::distance}).

\paragraph{Kiedy lista jest rzeczywiście potrzebna?}
Lista bywa przydatna, gdy musimy często \emph{przenosić węzły bez
kopiowania obiektów} i wymagamy stabilnych iteratorów. Dla częstej
iteracji i obliczeń numerycznych \texttt{vector} z reguły zapewnia
lepszą lokalność pamięci. Wybór należy jednak potwierdzić pomiarem.

\section{Kontenery uporządkowane: \texttt{set} i \texttt{map}}
\label{sec:cpp13-ordered}

\texttt{std::set<Key>} przechowuje unikalne klucze, a
\texttt{std::map<Key,T>} unikalne pary klucz--wartość. Elementy są
iterowane w kolejności zdefiniowanej przez komparator; wyszukiwanie,
wstawianie i usuwanie są logarytmiczne. Warianty
\texttt{multiset}/\texttt{multimap} dopuszczają powtarzające się klucze.

\subsection{\texttt{map::operator[]}, \texttt{at} i \texttt{find}}

Odwołanie \texttt{m[key]} do nieistniejącego klucza \emph{wstawia} nowy
element z wartością inicjalizowaną domyślnie; może zatem modyfikować
strukturę danych. \texttt{at()} zgłasza wyjątek, jeżeli klucz nie
istnieje. \texttt{find()} zwraca iterator, a \texttt{contains()} (C++20)
zwraca wartość logiczną bez wstawiania.

\begin{lstlisting}[caption={Liczenie czestosci slow z map (C++20); kompletny program.}]
#include <cassert>
#include <map>
#include <string>
#include <vector>

int main() {
    std::vector<std::string> words{"kot", "pies", "kot"};
    std::map<std::string, int> count;
    for (const auto& word : words) ++count[word];

    assert(count.at("kot") == 2);
    assert(count.contains("pies"));
    assert(!count.contains("lis"));
    assert(count.find("lis") == count.end());
}
\end{lstlisting}

\subsection{\texttt{insert}, \texttt{try\_emplace}, \texttt{insert\_or\_assign}}

\texttt{insert} podejmuje próbę dodania gotowej pary.
\texttt{try\_emplace(key,args...)} (C++17) konstruuje wartość
w miejscu z przekazanych argumentów tylko wtedy, gdy klucz nie występuje.
Ważne zastrzeżenie: \emph{wyrażenia będące argumentami wywołania}
są obliczane jeszcze przed wejściem do \texttt{try\_emplace}; nie
należy więc podawać niepotrzebnie kosztownego wyrażenia tworzącego
tymczasowy obiekt. \texttt{insert\_or\_assign} dodaje lub zastępuje wartość.

\begin{lstlisting}[caption={Rozne operacje modyfikacji map (C++17); kompletny program.}]
#include <cassert>
#include <map>
#include <memory>
#include <string>

struct Heavy {
    static inline int constructions = 0; // C++17
    int value;
    explicit Heavy(int x) : value(x) { ++constructions; }
};

int main() {
    std::map<std::string, Heavy> m;
    auto [it, added] = m.try_emplace("a", 3);
    assert(added && it->second.value == 3);

    auto [again, added_again] = m.try_emplace("a", 4);
    assert(!added_again && again->second.value == 3);
    assert(Heavy::constructions == 1);

    m.insert_or_assign("a", Heavy{9});
    assert(m.at("a").value == 9);

    std::map<int, std::unique_ptr<int>> owned;
    owned.try_emplace(1, std::make_unique<int>(10));
    assert(*owned.at(1) == 10);
}
\end{lstlisting}

\subsection{Własny komparator i ścisły słaby porządek}

Porównanie musi spełniać warunek \emph{strict weak ordering}; w
szczególności komparator powinien być irrefleksyjny
($\neg\operatorname{comp}(x,x)$) i przechodni. Jeżeli dwie wartości
$ a,b $ nie spełniają ani \texttt{comp(a,b)}, ani \texttt{comp(b,a)},
kontener uznaje je za równoważne jako klucze --- nawet gdy
\texttt{operator==} zwraca \texttt{false}. Wymagania tego porządku
są istotne również dla \texttt{sort} i \texttt{lower\_bound}.
Na przykład komparator ignorujący wielkość liter uzna \texttt{"ABC"}
i \texttt{"abc"} za ten sam klucz: druga próba wstawienia do
\texttt{map} nie utworzy nowego elementu.

\begin{lstlisting}[caption={Mapowanie bez rozrozniania wielkosci liter: prosty komparator; kompletny program.}]
#include <cassert>
#include <cctype>
#include <map>
#include <string>

struct CaseInsensitiveLess {
    bool operator()(const std::string& a,
                    const std::string& b) const {
        auto ai = a.begin();
        auto bi = b.begin();
        while (ai != a.end() && bi != b.end()) {
            auto x = static_cast<unsigned char>(*ai++);
            auto y = static_cast<unsigned char>(*bi++);
            int lx = std::tolower(x);
            int ly = std::tolower(y);
            if (lx != ly) return lx < ly;
        }
        return a.size() < b.size();
    }
};

int main() {
    std::map<std::string, int, CaseInsensitiveLess> m;
    m.emplace("ABC", 1);
    m.emplace("abc", 2); // rownowazny klucz
    assert(m.size() == 1 && m.begin()->second == 1);
}
\end{lstlisting}

Przykład porównuje bajty w podstawowym środowisku znakowym;
\texttt{std::tolower} z bieżącą lokalizacją C nie realizuje ogólnego
Unicode case folding. Do pełnego porównywania napisów w Unicode trzeba
zastosować osobne, świadomie wybrane reguły normalizacji i kolacji.

\subsection{Wyszukiwanie heterogeniczne}

Jeżeli komparator mapy jest \emph{transparentny} (np. \texttt{std::less<>}),
możemy szukać w mapie o kluczach \texttt{std::string} za pomocą
\texttt{std::string\_view}, bez tworzenia tymczasowego ciągu znaków.
Heterogeniczne przeciążenia \texttt{find} są dostępne
od C++14, a \texttt{contains} od C++20.

\begin{lstlisting}[caption={Heterogeniczne wyszukiwanie w map (C++20); kompletny program.}]
#include <cassert>
#include <functional>
#include <map>
#include <string>
#include <string_view>

int main() {
    std::map<std::string, int, std::less<>> m{{"alpha", 3}};
    std::string_view key = "alpha";
    auto it = m.find(key);
    assert(it != m.end() && it->second == 3);
    assert(m.contains(key));
}
\end{lstlisting}

\subsection{Zmiana klucza przez uchwyt węzła (C++17)}

Kluczy w \texttt{map} nie wolno modyfikować bezpośrednio przez iterator.
Uchwyt \texttt{node\_handle} uzyskany przez \texttt{extract()} pozwala
wyjąć węzeł, zmienić jego klucz i wstawić z powrotem. Trzeba pamiętać,
że ponowne wstawienie może się nie udać z powodu kolizji kluczy.

\begin{lstlisting}[caption={Zmiana klucza przez extract (C++17); kompletny program.}]
#include <cassert>
#include <map>
#include <string>
#include <utility>

int main() {
    std::map<std::string, int> m{{"old", 7}};
    auto node = m.extract("old");
    assert(!node.empty());
    node.key() = "new";
    auto result = m.insert(std::move(node));
    assert(result.inserted && m.at("new") == 7);
}
\end{lstlisting}

\section{Kontenery haszujące: \texttt{unordered\_map} i \texttt{unordered\_set}}
\label{sec:cpp13-unordered}

Kontener \texttt{unordered\_map<Key,T,Hash,Eq>} rozmieszcza elementy
w kubełkach. Funkcja \texttt{Hash} przekształca klucz w wartość typu
\texttt{std::size\_t}; \texttt{Eq} rozstrzyga równość kluczy.
Wymagany jest warunek zgodności:
\[
\operatorname{Eq}(a,b) \Longrightarrow
\operatorname{Hash}(a)=\operatorname{Hash}(b).
\]
Warunek odwrotny nie zachodzi: kolizje haszy są dozwolone i
rozstrzygane przez porównanie kluczy. Iteracja po kontenerze
nie ma stabilnego porządku i nie należy z niego korzystać jako z
interfejsu logicznego programu.
Można wyobrazić sobie hasz jako wskazówkę, w którym kubełku szukać.
Jeśli dwa różne klucze trafią do tego samego kubełka, \texttt{Eq}
sprawdza, czy są naprawdę równe. Stała w poniższym \texttt{PointHash}
służy jedynie do wymieszania haszy obu współrzędnych.

\begin{lstlisting}[caption={Wlasny hash i porownanie kluczy (C++20); kompletny program.}]
#include <cassert>
#include <cstddef>
#include <functional>
#include <unordered_map>

struct Point {
    int x, y;
    friend bool operator==(const Point&, const Point&) = default;
};

struct PointHash {
    std::size_t operator()(const Point& p) const noexcept {
        std::size_t h1 = std::hash<int>{}(p.x);
        std::size_t h2 = std::hash<int>{}(p.y);
        return h1 ^ (h2 + std::size_t{0x9e3779b9}
                        + (h1 << 6) + (h1 >> 2));
    }
};

int main() {
    std::unordered_map<Point, int, PointHash> m;
    m[{2, 5}] = 42;
    assert(m.at(Point{2, 5}) == 42);
}
\end{lstlisting}

Ta funkcja mieszająca służy demonstracji składni, a nie zastosowaniom
kryptograficznym. Dla publicznie sterowanych kluczy przy projektowaniu
systemów wystawionych na ataki kolizyjne należy rozważyć również
odporność na złośliwe dane wejściowe.

\subsection{\texttt{load\_factor}, \texttt{reserve} i \texttt{rehash}}

\texttt{load\_factor()} to stosunek liczby elementów do liczby kubełków.
\texttt{max\_load\_factor()} podaje lub ustawia pożądany maksymalny
współczynnik. \texttt{reserve(n)} żąda przygotowania na $n$ elementów,
a \texttt{rehash(k)} na przynajmniej $k$ kubełków, z uwzględnieniem
limitów wynikających z współczynnika wypełnienia.

\begin{lstlisting}[caption={Rezerwowanie kubelek przed licznymi wstawkami (C++20); kompletny program.}]
#include <cassert>
#include <unordered_set>

int main() {
    std::unordered_set<int> seen;
    seen.reserve(1000);
    for (int i = 0; i < 1000; ++i) seen.insert(i);
    assert(seen.size() == 1000);
    assert(seen.contains(17));
}
\end{lstlisting}

Rehashing unieważnia iteratory, lecz \emph{nie unieważnia} wskaźników
i referencji do elementów. Wstawienie może unieważnić iteratory, jeśli
wywoła rehash. Usunięcie unieważnia iteratory i referencje do usuniętych
elementów. Nie wolno modyfikować kluczy w sposób zaburzający ich hash
lub równość podczas przebywania elementów w kontenerze.

\section{Adaptery: \texttt{stack}, \texttt{queue}, \texttt{priority\_queue}}
\label{sec:cpp13-adaptors}

\texttt{stack} udostępnia schemat LIFO, \texttt{queue} --- FIFO,
a \texttt{priority\_queue} --- dostęp do elementu o najwyższym
priorytecie według komparatora (domyślnie do największego).
Adaptery zwykle nie zapewniają bezpośredniej iteracji po elementach.

\begin{lstlisting}[caption={Trzy adaptery standardowe (C++11); kompletny program.}]
#include <cassert>
#include <functional>
#include <queue>
#include <stack>
#include <vector>

int main() {
    std::stack<int> s;
    s.push(1); s.push(2);
    assert(s.top() == 2);
    s.pop();

    std::queue<int> q;
    q.push(1); q.push(2);
    assert(q.front() == 1);
    q.pop();

    std::priority_queue<int, std::vector<int>, std::greater<int>> pq;
    pq.push(7); pq.push(2); pq.push(5);
    assert(pq.top() == 2); // kopiec minimum
}
\end{lstlisting}

\texttt{priority\_queue} zwykle opiera się na kopcu: \texttt{top()} ma
koszt $O(1)$, a \texttt{push()} i \texttt{pop()} $O(\log n)$. Nie jest
to kolejka FIFO ani struktura z aktualizacją priorytetu dowolnego elementu
przez gotowy publiczny interfejs.

\section{Iteratory, stabilność adresów i unieważnianie}
\label{sec:cpp13-invalidation}

Iterator nie jest ogólnie tożsamy ze wskaźnikiem. Może przechowywać
stan niezbędny do poruszania się po drzewie lub kubełkach. Po operacji
modyfikującej trzeba wiedzieć, czy iterator, referencja i wskaźnik do
elementu pozostają ważne. Typowe reguły podsumowano w tabeli.

\begin{center}
\footnotesize
\begin{tabularx}{\textwidth}{@{}lXX@{}}
\toprule
\textbf{Kontener} & \textbf{Wstawienie} & \textbf{Usunięcie} \\
\midrule
\texttt{vector} & Przy realokacji: wszystkie iteratory i referencje; bez realokacji: pozycja wstawienia i dalej & Pozycja usunięcia i dalej (także dawny \texttt{end()}) \\
\texttt{deque} & Na końcach: iteratory tracą ważność, referencje do elementów zwykle pozostają; w środku: także referencje & Zależy od pozycji; usunięcie w środku unieważnia wszystkie iteratory i referencje \\
\texttt{list} & Istniejące iteratory i referencje zachowane & Tylko usunięte elementy \\
\texttt{map}, \texttt{set} & Istniejące iteratory i referencje zachowane & Tylko usunięte elementy \\
\texttt{unordered\_*} & Iteratory mogą utracić ważność przy rehash; referencje zachowane & Tylko usunięte elementy \\
\bottomrule
\end{tabularx}
\end{center}

Dla \texttt{deque} szczególnie ostrożnie trzeba traktować
\texttt{end()}, którego ważność zależy od rodzaju usunięcia.
W tabeli podano reguły najważniejsze dla praktyki, a nie kompletną
specyfikację każdej przeciążonej operacji. Po operacjach zmieniających
strukturę bezpieczną praktyką jest korzystanie z iteratora zwróconego
przez metodę (np. \texttt{erase}, \texttt{insert}).

\subsection{Bezpieczne usuwanie podczas iteracji}

\begin{lstlisting}[caption={Usuwanie elementow map podczas iteracji (C++20); kompletny program.}]
#include <cassert>
#include <map>

int main() {
    std::map<int, int> m{{1, 2}, {2, 3}, {3, 4}};
    for (auto it = m.begin(); it != m.end();) {
        if (it->second % 2 == 0) it = m.erase(it);
        else ++it;
    }
    assert(m.size() == 1 && m.contains(2));
}
\end{lstlisting}

\paragraph{Szczególny przypadek \texttt{vector<bool>}.}
\texttt{std::vector<bool>} jest specjalizacją oszczędzającą pamięć,
która może reprezentować wartości w postaci bitów. Jej \texttt{operator[]}
nie musi zwracać \texttt{bool\&}, tylko specjalny obiekt proxy. Nie
należy oczekiwać normalnej tablicy obiektów \texttt{bool}, ciągłego
przechowywania pojedynczych \texttt{bool} ani bezpiecznej modyfikacji
odrębnych bitów z różnych wątków bez synchronizacji.

\section{Algorytmy standardowe i wymagania wobec iteratorów}

\begin{center}\small
\begin{tabularx}{\textwidth}{@{}p{0.34\textwidth}X@{}}
\toprule
Algorytm & Działanie i warunki \\
\midrule
\code{std::find} & Szuka pierwszej wartości równej wzorcowi; przy braku zwraca iterator końcowy. \\
\code{std::transform} & Stosuje funkcję do zakresu i zapisuje wyniki przez iterator wyjściowy. \\
\code{std::sort} & Sortuje zakres z iteratorami dostępu swobodnego; nie zachowuje kolejności elementów równoważnych. \\
\code{std::stable_sort} & Sortuje, zachowując wzajemną kolejność elementów równoważnych. \\
\code{std::lower_bound} & Szuka pierwszego elementu nie mniejszego od wzorca; wymaga odpowiednio uporządkowanego zakresu. \\
\code{std::erase_if} & Usuwa z kontenera elementy spełniające predykat; dostępne od C++20 dla wspieranych kontenerów. \\
\bottomrule
\end{tabularx}
\end{center}
\label{sec:cpp13-algorithms}

Algorytmy oddzielają \emph{operację} od \emph{reprezentacji danych}.
Wersje klasyczne przyjmują iteratory, a w C++20 wiele algorytmów
w \texttt{std::ranges} można wywołać bezpośrednio na zakresie. Szczegóły
kategorii iteratorów i widoków omawiamy w rozdziale~\ref{chap:ranges-views-cpp}; tutaj skupiamy
się na doborze algorytmu i jego kontrakcie.

\subsection{Wyszukiwanie, transformacje i agregacje}

Do częstych operacji należą \texttt{find}, \texttt{find\_if},
\texttt{count\_if}, \texttt{transform}, \texttt{accumulate},
\texttt{all\_of}, \texttt{any\_of}, \texttt{none\_of} i
\texttt{for\_each}. Algorytmy numeryczne znajdują się w nagłówku
\texttt{<numeric>}, pozostałe głównie w \texttt{<algorithm>}.

\begin{lstlisting}[caption={Algorytmy klasyczne i Ranges (C++20); kompletny program.}]
#include <algorithm>
#include <cassert>
#include <numeric>
#include <vector>

int main() {
    std::vector<int> values{1, 2, 3, 4, 5};
    auto even = [](int x) { return x % 2 == 0; };

    assert(std::ranges::count_if(values, even) == 2);
    assert(std::ranges::all_of(values,
        [](int x) { return x > 0; }));
    assert(std::accumulate(values.begin(), values.end(), 0) == 15);

    std::vector<int> squares(values.size());
    std::ranges::transform(values, squares.begin(),
                           [](int x) { return x * x; });
    assert(squares[2] == 9);
}
\end{lstlisting}

\subsection{Sortowanie i stabilność kolejności}

\texttt{std::sort} wymaga iteratorów dostępu swobodnego i ma
złożoność $O(n\log n)$ w najgorszym przypadku (w liczbie porównań).
\texttt{stable\_sort} zachowuje kolejność elementów, które komparator
uznaje za równoważne, zwykle kosztem dodatkowych zasobów; jego
szczegółowe gwarancje złożoności zależą od dostępnej pamięci.
\texttt{nth\_element} umieszcza wybrany element na pozycji, którą
zajmowałby po sortowaniu, bez sortowania całej sekwencji; przeciętna
liczba porównań jest liniowa.

\begin{lstlisting}[caption={Stabilne sortowanie rekordow po kluczu (C++20); kompletny program.}]
#include <algorithm>
#include <cassert>
#include <string>
#include <vector>

struct Record {
    std::string name;
    int group;
};

int main() {
    std::vector<Record> v{{"A", 2}, {"B", 1}, {"C", 2}};
    std::ranges::stable_sort(v, {}, &Record::group);
    assert(v[0].name == "B");
    assert(v[1].name == "A" && v[2].name == "C");
}
\end{lstlisting}

\subsection{Wyszukiwanie binarne i porządek wejścia}

\texttt{binary\_search}, \texttt{lower\_bound}, \texttt{upper\_bound}
i \texttt{equal\_range} wymagają odpowiedniego uporządkowania
(lub partycjonowania względem danego klucza) zakresu. Zastosowanie
do nieuporządkowanych danych daje wynik niespełniający kontraktu
algorytmu. W wersjach na iteratorach nie-losowego dostępu liczba
porównań może pozostać logarytmiczna, lecz przesuwanie iteratorów
może być liniowe.

\begin{lstlisting}[caption={Wyszukiwanie miejsca wstawienia w posortowanym vector (C++11); kompletny program.}]
#include <algorithm>
#include <cassert>
#include <vector>

int main() {
    std::vector<int> sorted{1, 3, 3, 7};
    auto it = std::lower_bound(sorted.begin(), sorted.end(), 3);
    assert(it == sorted.begin() + 1);

    auto upper = std::upper_bound(sorted.begin(), sorted.end(), 3);
    assert(upper == sorted.begin() + 3);
    assert(std::binary_search(sorted.begin(), sorted.end(), 7));
}
\end{lstlisting}

\subsection{Idiom \emph{erase--remove} i \texttt{std::erase\_if}}

\texttt{std::remove} nie usuwa elementów \emph{z kontenera}:
przestawia elementy pozostające i zwraca logiczny nowy koniec.
Rozmiar kontenera zmniejsza dopiero wywołanie \texttt{erase}.
Od C++20 wygodniej używać \texttt{std::erase} i
\texttt{std::erase\_if} dla wspieranych kontenerów.

\begin{lstlisting}[caption={Dwa sposoby usuwania z vector (C++20); kompletny program.}]
#include <algorithm>
#include <cassert>
#include <vector>

int main() {
    std::vector<int> a{1, 2, 3, 4, 5};
    auto last = std::remove_if(a.begin(), a.end(),
                               [](int x) { return x % 2 == 0; });
    a.erase(last, a.end());
    assert((a == std::vector<int>{1, 3, 5}));

    std::vector<int> b{1, 2, 3, 4, 5};
    auto removed = std::erase_if(b,
                                 [](int x) { return x % 2 == 0; });
    assert(removed == 2 && b == a);
}
\end{lstlisting}

Dla kontenerów asocjacyjnych nie wolno stosować klasycznego
\texttt{remove\_if} tak jak dla wektora: ich elementy mają klucze,
których pozycja jest regulowana przez strukturę danych. Stosujemy
\texttt{erase} z iteratorem lub \texttt{std::erase\_if} (C++20).

\subsection{Algorytmy kopca i \texttt{std::priority\_queue}}

\texttt{make\_heap}, \texttt{push\_heap}, \texttt{pop\_heap} i
\texttt{sort\_heap} pozwalają użyć własnego wektora jako kopca.
\texttt{make\_heap} działa w czasie liniowym, a operacje dopisania i
usunięcia maksimum w $O(\log n)$. Są to operacje inne niż pełne sortowanie.

\section{Przypadki projektowe: dobór kontenera do problemu}
\label{sec:cpp13-use-cases}

\subsection{Histogram, deduplikacja i lista wyników}

Załóżmy, że chcemy zebrać wartości, usunąć duplikaty i otrzymać
wynik w kolejności rosnącej. Możemy użyć \texttt{std::set}, albo
\texttt{vector} z \texttt{sort} i \texttt{unique}. Drugi sposób bywa
korzystniejszy w przetwarzaniu wsadowym, ponieważ \texttt{vector}
przechowuje dane ciągle i ogranicza liczbę alokacji węzłów.

\begin{lstlisting}[caption={Deduplikacja przez sort i unique (C++11); kompletny program.}]
#include <algorithm>
#include <cassert>
#include <vector>

int main() {
    std::vector<int> keys{5, 2, 5, 1, 2, 8};
    std::sort(keys.begin(), keys.end());
    keys.erase(std::unique(keys.begin(), keys.end()), keys.end());
    assert((keys == std::vector<int>{1, 2, 5, 8}));
}
\end{lstlisting}

\subsection{Własność elementów w kontenerach polimorficznych}

Przechowywanie obiektów klas pochodnych przez wartość w
\texttt{vector<Base>} może powodować \emph{object slicing}. Dla
polimorfizmu dynamicznego używamy np. wektora wskaźników
\texttt{unique\_ptr<Base>}. Przy zamkniętej rodzinie typów można użyć
wektora z elementami \texttt{variant} (zob. rozdział~\ref{chap:cpp14-type-erasure}).
Różnicę między polimorfizmem a obcięciem obiektu omawia
rozdział~\ref{chap:cpp-object-model}.

\begin{lstlisting}[caption={Kontener wlasciwych obiektow polimorficznych (C++14); kompletny program.}]
#include <cassert>
#include <memory>
#include <vector>

struct Shape {
    virtual ~Shape() = default;
    virtual int area() const = 0;
};
struct Square : Shape {
    int side;
    explicit Square(int x) : side(x) {}
    int area() const override { return side * side; }
};

int main() {
    std::vector<std::unique_ptr<Shape>> shapes;
    shapes.push_back(std::make_unique<Square>(3));
    assert(shapes.front()->area() == 9);
}
\end{lstlisting}

\subsection{Kontenery a współbieżność}

Samo wybranie kontenera standardowego nie czyni jego operacji
modyfikujących bezpiecznymi dla współbieżnego wykonania. Równoczesne
modyfikowanie wspólnego \texttt{vector} przez niezależne
\texttt{push\_back()} bez synchronizacji jest wyścigiem danych.
Można natomiast równolegle odczytywać elementy kontenera, jeśli żaden
wątek ich nie modyfikuje, a czas życia kontenera jest zachowany.
Dla wielu kontenerów można także modyfikować \emph{różne istniejące}
elementy równolegle, jeśli nie dochodzi do konfliktu pamięci; wyjątek
stanowi m.in. specjalizacja \texttt{vector<bool>}. Zmiany strukturalne
wymagają osobnej analizy i synchronizacji (zob. rozdział~\ref{chap:concurrency}).

\section{Złożoność, lokalność pamięci i pomiar wydajności}
\label{sec:cpp13-performance}

Teoretyczna złożoność to jeden z elementów oceny. Istotne są również:
liczba alokacji, liczba odwołań do rozproszonych adresów, rozmiar
elementu, przewidywalność gałęzi i zachowanie cache procesora.
Dla wielu zastosowań \texttt{vector} z wyszukiwaniem liniowym
może być szybszy od struktury drzewiastej przy małej liczbie kluczy.
To nie przeczy złożoności asymptotycznej: stałe i charakter obciążenia
są częścią kosztu rzeczywistego programu.

\subsection{Koszt sumowania ciągłych danych}

\begin{lstlisting}[caption={Prosty szkielet pomiaru vector i list (C++20); kompletny program.}]
#include <chrono>
#include <cstdint>
#include <iostream>
#include <list>
#include <numeric>
#include <utility>
#include <vector>

int main() {
    constexpr int N = 200000;
    std::vector<int> v(N, 1);
    std::list<int> l(v.begin(), v.end());

    auto measure = [](const auto& xs) {
        auto begin = std::chrono::steady_clock::now();
        std::int64_t sum = std::accumulate(xs.begin(), xs.end(),
                                           std::int64_t{0});
        auto end = std::chrono::steady_clock::now();
        auto ns = std::chrono::duration_cast<
            std::chrono::nanoseconds>(end - begin).count();
        return std::pair{sum, ns};
    };

    auto [sum_v, time_v] = measure(v);
    auto [sum_l, time_l] = measure(l);
    std::cout << "vector: " << sum_v << ", " << time_v << " ns\n";
    std::cout << "list:   " << sum_l << ", " << time_l << " ns\n";
}
\end{lstlisting}

To \emph{ilustracja sposobu pomiaru}, a nie miarodajny benchmark.
Jednorazowy wynik bywa zdominowany przez rozgrzanie cache, optymalizacje
kompilatora, scheduler i zakłócenia systemowe. Poprawniejszy eksperyment
powtarza pomiary, rozdziela przygotowanie danych od mierzonej operacji,
uwzględnia pamięć, różne rozmiary danych, medianę i rozrzut pomiarów.
Dla wiarygodnych porównań warto użyć np. Google Benchmark i profilera.

\subsection{Przybliżone reguły wyboru}

\begin{center}
\small
\begin{tabularx}{\textwidth}{@{}lX@{}}
\toprule
\textbf{Wymaganie} & \textbf{Rozsądny wybór początkowy} \\
\midrule
Dużo iteracji, obliczeń, dostęp indeksowany & \texttt{vector}; przy stałym rozmiarze \texttt{array} \\
Częste dodawanie i usuwanie z obu końców & \texttt{deque} \\
Utrzymanie kolejności kluczy, wyszukiwanie przedziałowe & \texttt{map} / \texttt{set} \\
Wyszukiwanie po kluczu bez wymogu kolejności & \texttt{unordered\_map} / \texttt{unordered\_set} \\
Przenoszenie istniejących węzłów i stabilne iteratory & \texttt{list} \\
FIFO / LIFO & \texttt{queue} / \texttt{stack} \\
Wielokrotne pobieranie największego lub najmniejszego & \texttt{priority\_queue} \\
\bottomrule
\end{tabularx}
\end{center}

\section{Najczęstsze błędy i dobre praktyki}
\label{sec:cpp13-pitfalls}

\begin{enumerate}
\item Używanie \texttt{vector::operator[]} bez pewności, że indeks jest
  poprawny; w kodzie granicznym warto stosować \texttt{at()}.
\item Traktowanie \texttt{reserve} jako tworzenia elementów (mylenie
  \texttt{size} z \texttt{capacity}).
\item Korzystanie z iteratorów lub referencji po unieważnieniu przez
  realokację, \texttt{erase} lub \texttt{rehash}.
\item Zakładanie, że \texttt{emplace\_back} zawsze jest szybsze od
  \texttt{push\_back}.
\item Używanie \texttt{map::operator[]} do samego sprawdzania istnienia
  klucza, z niezamierzonym wstawieniem nowej wartości.
\item Pisanie komparatorów naruszających \emph{strict weak ordering};
  komparator \texttt{a <= b} nie jest poprawną relacją porządkującą.
\item Używanie \texttt{std::sort} na \texttt{std::list} lub
  \texttt{std::lower\_bound} na danych bez odpowiedniego porządku.
\item Wnioskowanie o kolejności elementów \texttt{unordered\_map}
  na podstawie przypadkowego przebiegu programu.
\item Niepotrzebne stosowanie \texttt{list} z przekonaniem, że
  wstawienie $O(1)$ zawsze wygrywa z \texttt{vector}.
\item Wprowadzanie równoczesnych modyfikacji kontenera z wielu wątków
  bez uzasadnienia braku konfliktu i odpowiedniej synchronizacji.
\end{enumerate}

\section{Pytania kontrolne}

\begin{enumerate}
\item Jaka jest różnica między \texttt{reserve} a \texttt{resize}, i
  które z nich tworzy nowe elementy?
\item Kiedy \texttt{push\_back} na \texttt{vector} unieważnia wszystkie
  referencje, a kiedy zachowuje referencje do wcześniejszych elementów?
\item Dlaczego \texttt{map::operator[]} może być niepożądany podczas
  sprawdzania istnienia klucza?
\item Jakie wymagania muszą spełniać \texttt{Hash} i \texttt{Eq} w
  \texttt{unordered\_map}?
\item Dlaczego \texttt{std::remove} nie zmienia rozmiaru \texttt{vector}?
\item Jak różnią się koszty i kontrakty \texttt{sort}, \texttt{stable\_sort},
  \texttt{nth\_element} i \texttt{lower\_bound}?
\item W jakich okolicznościach \texttt{list::splice} zapewnia korzyść
  niedostępną dla zwykłego \texttt{vector}?
\end{enumerate}

\section{Podsumowanie rozdziału}

Dobry kod C++ wybiera kontener zgodnie z potrzebnymi operacjami i
własnościami pamięci, a nie jedynie na podstawie pojedynczego rzędu
złożoności. \texttt{vector} to często najlepszy punkt wyjścia;
\texttt{deque}, kontenery drzewiaste, haszujące i listy rozwiązują
bardziej szczegółowe problemy. Praktycznie równie ważne jak złożoności
są reguły unieważniania iteratorów, wymagania komparatorów,
bezpieczeństwo czasu życia i świadome wykorzystywanie algorytmów
standardowych zamiast ręcznie pisanych pętli.

\section*{Literatura i dokumentacja}
\addcontentsline{toc}{section}{Literatura i dokumentacja}
\begin{itemize}
\item \url{https://en.cppreference.com/w/cpp/container} --- kontenery i
  zestawienie reguł unieważniania iteratorów.
\item \url{https://en.cppreference.com/w/cpp/algorithm} --- algorytmy
  biblioteki standardowej i ich kontrakty.
\item \url{https://en.cppreference.com/w/cpp/container/vector} ---
  \texttt{std::vector} i zarządzanie pojemnością.
\item \url{https://en.cppreference.com/w/cpp/container/unordered_map} ---
  struktury haszujące.
\item Scott Meyers, \emph{Effective Modern C++}, O'Reilly, 2014.
\item Nicolai M. Josuttis, \emph{The C++ Standard Library}, wydanie drugie,
  Addison-Wesley, 2012.
\end{itemize}


\chapter{Kontenery STL i algorytmy}
\label{chap:cpp-stl-containers}

\section{Cel i zakres rozdziału}

Kontener przechowuje grupę elementów i pozwala je odnajdywać lub
zmieniać. Najprostsza mapa wyboru wygląda tak: \texttt{vector} jest
listą z dostępem przez indeks, \texttt{map} słownikiem klucz--wartość,
a \texttt{set} zbiorem wartości bez powtórzeń. Jeśli kolejność kluczy
nie ma znaczenia, istnieją także haszujące warianty
\texttt{unordered\_map} i \texttt{unordered\_set}.

Biblioteka standardowa C++ udostępnia gotowe struktury danych, których wybór ma
istotne konsekwencje dla złożoności operacji, lokalności pamięci, bezpieczeństwa
wyjątkowego i poprawności iteratorów. Samo rozróżnienie między
\texttt{std::vector}, \texttt{std::list} i \texttt{std::map} nie wystarcza:
trzeba rozumieć również \emph{własność elementów}, \emph{unieważnianie
iteratorów}, wymagania dotyczące porządku i haszowania oraz relację między
kontenerami a algorytmami. W dalszej części domyślnie używamy C++20.

\paragraph{Główna zasada projektowa.}
Najpierw ustalamy rzeczywiste operacje: czy potrzebujemy dostępu przez
indeks, wyszukiwania po kluczu, zachowania kolejności, stabilnych adresów,
czy częstych wstawek i usunięć. Dopiero potem wybieramy kontener. Prosty
\texttt{std::vector} pozostaje często najlepszym rozwiązaniem, nawet gdy
jego teoretyczna złożoność niektórych operacji jest gorsza, ponieważ dane
leżą w sąsiadujących komórkach pamięci.

\section{Klasyfikacja kontenerów i podstawowe pojęcia}
\label{sec:cpp13-taxonomy}

Biblioteka standardowa rozróżnia m.in. kontenery sekwencyjne
(\texttt{array}, \texttt{vector}, \texttt{deque}, \texttt{list},
\texttt{forward\_list}), uporządkowane asocjacyjne (\texttt{map},
\texttt{set} i ich warianty \texttt{multi}) oraz nieuporządkowane
asocjacyjne (\texttt{unordered\_map}, \texttt{unordered\_set} i warianty
\texttt{multi}). \texttt{stack}, \texttt{queue} i \texttt{priority\_queue}
są \emph{adapterami kontenerów}: oferują inny interfejs dostępu do
wewnętrznej struktury danych. \texttt{std::string} jest wyspecjalizowaną
strukturą sekwencyjną, ale semantykę tekstu omawiamy oddzielnie.

\begin{center}
\small
\begin{tabularx}{\textwidth}{@{}lXXX@{}}
\toprule
\textbf{Kontener} & \textbf{Dostęp} & \textbf{Wstawienie / usunięcie} & \textbf{Typowe zastosowanie} \\
\midrule
\texttt{vector} & Indeks $O(1)$ & Koniec: zamortyzowane $O(1)$; środek: $O(n)$ & Dane zwarte, iteracja, obliczenia \\
\texttt{deque} & Indeks $O(1)$ & Oba końce: $O(1)$; środek: $O(n)$ & Kolejka dwustronna \\
\texttt{list} & Iteracja $O(n)$ & $O(1)$ przy znanym iteratorze & Stabilne węzły, \texttt{splice} \\
\texttt{forward\_list} & Iteracja $O(n)$ & $O(1)$ za poprzednikiem & Lista jednokierunkowa \\
\texttt{map} / \texttt{set} & Szukanie $O(\log n)$ & $O(\log n)$ & Klucze w ustalonej kolejności \\
\texttt{unordered\_map} & Średnio $O(1)$ & Średnio $O(1)$ & Słownik, szybkie wyszukiwanie \\
\bottomrule
\end{tabularx}
\end{center}

W tabeli $n$ oznacza liczbę elementów, a koszty zależą od wymagań
poszczególnych operacji. \texttt{unordered\_map} może w najgorszym
przypadku wykonywać operacje w czasie $O(n)$ (np. przy wielu kolizjach).
Sformułowanie ,,\emph{zamortyzowane} $O(1)$'' oznacza, że pojedyncza
operacja może być kosztowna, ale średni koszt w długiej sekwencji operacji
jest ograniczony stałą.

\paragraph{Wartości, wskaźniki, referencje.}
Kontenery zwykle przechowują \emph{własne obiekty elementów}.
\texttt{std::vector<T>} nie jest kontenerem referencji do obiektów
\texttt{T}; może za to przechowywać np. \texttt{std::reference\_wrapper<T>},
\texttt{T*} lub \texttt{std::unique\_ptr<T>}. Wybór wpływa na własność
obiektów, kopiowanie, czas życia i poprawność dostępu.

\section{\texttt{std::vector}: domyślny kontener sekwencyjny}

\begin{center}\small
\begin{tabularx}{\textwidth}{@{}p{0.34\textwidth}X@{}}
\toprule
Operacja & Za co odpowiada \\
\midrule
\code{size()}, \code{empty()} & Podają liczbę elementów i informację, czy kontener jest pusty. \\
\code{capacity()} & Podaje liczbę elementów, które zmieszczą się bez ponownej alokacji. \\
\code{reserve(n)} & Rezerwuje miejsce na co najmniej \code{n} elementów; nie zmienia \code{size()}. \\
\code{resize(n)} & Zmienia liczbę elementów: tworzy nowe albo niszczy nadmiarowe. \\
\code{push_back(x)} & Dodaje istniejącą wartość na koniec przez kopiowanie lub przeniesienie. \\
\code{emplace_back(args...)} & Konstruuje nowy element na końcu z podanych argumentów. \\
\code{erase(pos)} & Usuwa element; kolejne elementy przesuwają się, a iteratory od miejsca usunięcia tracą ważność. \\
\bottomrule
\end{tabularx}
\end{center}
\label{sec:cpp13-vector}

\texttt{std::vector<T>} przechowuje elementy \emph{ciągle} w pamięci,
co pozwala przekazywać je do interfejsów oczekujących tablicy (np. przez
\texttt{data()} lub \texttt{std::span}). Obiekt wektora zarządza własną
pamięcią i niszczy elementy po zakończeniu swojego czasu życia.

\begin{lstlisting}[caption={Podstawowe operacje na vector (C++11); kompletny program.}]
#include <cassert>
#include <vector>

int main() {
    std::vector<int> v{4, 8, 15};
    v.push_back(16);           // kopiowanie lub przenoszenie argumentu
    v.emplace_back(23);        // konstruowanie nowego elementu

    assert(v.size() == 5);
    assert(v[0] == 4);
    assert(v.at(4) == 23);
    assert(v.front() == 4 && v.back() == 23);

    v.pop_back();
    assert(v.size() == 4);
    v.clear();
    assert(v.empty());
}
\end{lstlisting}

\texttt{operator[]} \emph{nie sprawdza} granic: użycie indeksu spoza zakresu
prowadzi do niezdefiniowanego zachowania. Metoda \texttt{at()} sprawdza
zakres i zgłasza \texttt{std::out\_of\_range}. Metody \texttt{front()},
\texttt{back()} i \texttt{pop\_back()} wymagają niepustego kontenera.

\subsection{\texttt{size}, \texttt{capacity}, \texttt{reserve} i \texttt{resize}}

\texttt{size()} oznacza liczbę \emph{istniejących} elementów, a
\texttt{capacity()} liczbę elementów, dla których przydzielono miejsce
bez konieczności realokacji. \texttt{reserve(k)} zmienia rezerwację
pamięci, nie tworzy elementów; \texttt{resize(k)} zmienia ich liczbę.

\begin{lstlisting}[caption={Rezerwacja pamięci nie zwiększa rozmiaru (C++11); kompletny program.}]
#include <cassert>
#include <vector>

int main() {
    std::vector<int> data;
    data.reserve(100);
    assert(data.size() == 0);
    assert(data.capacity() >= 100);

    data.push_back(7);
    data.resize(4);          // dopisuje trzy elementy o wartosci 0
    assert(data.size() == 4);
    assert(data[0] == 7 && data[3] == 0);
}
\end{lstlisting}

Nie należy wywoływać \texttt{reserve(size()+1)} przed każdym
\texttt{push\_back()}: wymuszenie niemal liniowego wzrostu pojemności
może zniweczyć politykę geometrycznego powiększania bufora i pogorszyć
wydajność. \texttt{shrink\_to\_fit()} jest tylko żądaniem, a nie gwarancją
zmniejszenia pojemności.

\subsection{Realokacja i unieważnianie iteratorów}

Po realokacji \texttt{vector} zmienia adresy wszystkich elementów:
tracą ważność iteratory, referencje i wskaźniki do elementów, również
poprzednia wartość \texttt{data()}. Bez realokacji \texttt{push\_back()}
unieważnia iterator \texttt{end()}, lecz nie unieważnia iteratorów i
referencji do wcześniej istniejących elementów. Wstawienie w środku
bez realokacji unieważnia iteratory i referencje do pozycji wstawienia
oraz wszystkich późniejszych pozycji (włącznie z dawnym \texttt{end()}).

\begin{lstlisting}[caption={Bezpieczne ponowne pobranie iteratora po wstawieniu; kompletny program.}]
#include <cassert>
#include <vector>

int main() {
    std::vector<int> v{1, 3, 4};
    auto pos = v.begin() + 1;
    pos = v.insert(pos, 2);   // przyjmuje zwrocony iterator
    assert(*pos == 2);
    assert((v == std::vector<int>{1, 2, 3, 4}));
}
\end{lstlisting}

\paragraph{Przenoszenie elementów i bezpieczeństwo wyjątkowe.}
Podczas realokacji \texttt{vector} może przenosić lub kopiować elementy.
Dla pewnych typów i operacji warunki silnej gwarancji wyjątkowej zależą
od tego, czy konstruktor przenoszący jest \texttt{noexcept} oraz czy
dostępne jest kopiowanie. Zagadnienie łączy się bezpośrednio z semantyką
przenoszenia omówioną w rozdziale 6.

\subsection{\texttt{push\_back} kontra \texttt{emplace\_back}}

\texttt{emplace\_back(args...)} przekazuje argumenty bezpośrednio do
konstruktora nowego elementu. Nie oznacza to jednak automatycznie
,,zawsze szybciej'': jeśli mamy już obiekt, \texttt{push\_back(std::move(x))}
bywa równie naturalne; przy realokacji nadal konieczne może być
przeniesienie innych elementów. Obie funkcje mogą zgłaszać wyjątki.

\begin{lstlisting}[caption={Konstruowanie elementu z argumentow (C++17); kompletny program.}]
#include <cassert>
#include <string>
#include <utility>
#include <vector>

struct Person {
    std::string name;
    int age;
    Person(std::string n, int a) : name(std::move(n)), age(a) {}
};

int main() {
    std::vector<Person> people;
    people.emplace_back("Ada", 31);
    people.push_back(Person{"Alan", 42});
    assert(people[1].age == 42);
}
\end{lstlisting}

\section{\texttt{std::array}, \texttt{std::deque} i listy}
\label{sec:cpp13-sequential}

\subsection{\texttt{std::array<T,N>}: rozmiar znany w kompilacji}

\texttt{std::array} przechowuje dokładnie $N$ elementów. Nie dokonuje
osobnej dynamicznej alokacji dla swojej zawartości; pamięć elementów
stanowi część samego obiektu \texttt{array}. Można używać iteratorów,
\texttt{size()}, \texttt{data()} i algorytmów, analogicznie do wektora.

\begin{lstlisting}[caption={std::array i algorytmy standardowe (C++17); kompletny program.}]
#include <algorithm>
#include <array>
#include <cassert>
#include <tuple>

int main() {
    std::array<int, 4> x{5, 1, 4, 2};
    std::sort(x.begin(), x.end());
    assert((x == std::array<int, 4>{1, 2, 4, 5}));
    static_assert(std::tuple_size_v<decltype(x)> == 4);
}
\end{lstlisting}

\subsection{\texttt{std::deque}: dwa tanie końce}

\texttt{deque} udostępnia dostęp indeksowany w czasie $O(1)$, lecz
zwykle nie jest jednym ciągłym blokiem pamięci. Wstawianie i usuwanie
na początku i końcu ma stałą złożoność; w przeciwieństwie do
\texttt{vector} ma operacje \texttt{push\_front()} i \texttt{pop\_front()}.
\texttt{deque} nie udostępnia \texttt{data()} dla wszystkich elementów.

\begin{lstlisting}[caption={Kolejka dwustronna (C++11); kompletny program.}]
#include <cassert>
#include <deque>

int main() {
    std::deque<int> q{2, 3};
    q.push_front(1);
    q.push_back(4);
    assert(q[0] == 1 && q[3] == 4);
    q.pop_front();
    assert(q.front() == 2);
}
\end{lstlisting}

Wstawienie na końcu \texttt{deque} może unieważnić \emph{iteratory},
ale zachowuje referencje do istniejących elementów. Szczegółowe reguły
różnią się dla wstawień i usunięć w środku oraz na końcach; nie należy
mechanicznie przenosić reguł \texttt{vector} na \texttt{deque}.

\subsection{\texttt{std::list} i \texttt{std::forward\_list}}

\texttt{list} jest listą dwukierunkową. Wstawianie i usuwanie przy
znanym iteratorze odbywa się w stałym czasie; iteratory i referencje
do pozostałych elementów pozostają ważne. Wyszukanie elementu przez
wartość wymaga jednak przejścia listy $O(n)$, podobnie jak uzyskanie
elementu o zadanym indeksie. W przeciwieństwie do \texttt{vector},
\texttt{list} nie oferuje operatora \texttt{[]} ani iteratorów dostępu
swobodnego. Do sortowania używamy \texttt{list::sort()}, a nie
\texttt{std::sort()}.

\begin{lstlisting}[caption={Przeniesienie wezla przez list::splice (C++11); kompletny program.}]
#include <cassert>
#include <iterator>
#include <list>

int main() {
    std::list<int> source{1, 2, 3};
    std::list<int> target{10};
    auto it = std::next(source.begin()); // wskazuje 2

    target.splice(target.end(), source, it);
    assert(*it == 2); // iterator pozostaje wazny
    assert(source.size() == 2 && target.size() == 2);
    target.sort();
    assert(target.front() == 2);
}
\end{lstlisting}

\texttt{forward\_list} przechowuje węzły listy jednokierunkowej;
wstawianie i usuwanie wykonuje \emph{za} wskazanym poprzednikiem
(\texttt{insert\_after}, \texttt{erase\_after}), a \texttt{before\_begin()}
udostępnia pozycję przed pierwszym elementem. Kontener nie posiada
metody \texttt{size()} (rozmiar można policzyć liniowo przez
\texttt{std::distance}).

\paragraph{Kiedy lista jest rzeczywiście potrzebna?}
Lista bywa przydatna, gdy musimy często \emph{przenosić węzły bez
kopiowania obiektów} i wymagamy stabilnych iteratorów. Dla częstej
iteracji i obliczeń numerycznych \texttt{vector} z reguły zapewnia
lepszą lokalność pamięci. Wybór należy jednak potwierdzić pomiarem.

\section{Kontenery uporządkowane: \texttt{set} i \texttt{map}}
\label{sec:cpp13-ordered}

\texttt{std::set<Key>} przechowuje unikalne klucze, a
\texttt{std::map<Key,T>} unikalne pary klucz--wartość. Elementy są
iterowane w kolejności zdefiniowanej przez komparator; wyszukiwanie,
wstawianie i usuwanie są logarytmiczne. Warianty
\texttt{multiset}/\texttt{multimap} dopuszczają powtarzające się klucze.

\subsection{\texttt{map::operator[]}, \texttt{at} i \texttt{find}}

Odwołanie \texttt{m[key]} do nieistniejącego klucza \emph{wstawia} nowy
element z wartością inicjalizowaną domyślnie; może zatem modyfikować
strukturę danych. \texttt{at()} zgłasza wyjątek, jeżeli klucz nie
istnieje. \texttt{find()} zwraca iterator, a \texttt{contains()} (C++20)
zwraca wartość logiczną bez wstawiania.

\begin{lstlisting}[caption={Liczenie czestosci slow z map (C++20); kompletny program.}]
#include <cassert>
#include <map>
#include <string>
#include <vector>

int main() {
    std::vector<std::string> words{"kot", "pies", "kot"};
    std::map<std::string, int> count;
    for (const auto& word : words) ++count[word];

    assert(count.at("kot") == 2);
    assert(count.contains("pies"));
    assert(!count.contains("lis"));
    assert(count.find("lis") == count.end());
}
\end{lstlisting}

\subsection{\texttt{insert}, \texttt{try\_emplace}, \texttt{insert\_or\_assign}}

\texttt{insert} podejmuje próbę dodania gotowej pary.
\texttt{try\_emplace(key,args...)} (C++17) konstruuje wartość
w miejscu z przekazanych argumentów tylko wtedy, gdy klucz nie występuje.
Ważne zastrzeżenie: \emph{wyrażenia będące argumentami wywołania}
są obliczane jeszcze przed wejściem do \texttt{try\_emplace}; nie
należy więc podawać niepotrzebnie kosztownego wyrażenia tworzącego
tymczasowy obiekt. \texttt{insert\_or\_assign} dodaje lub zastępuje wartość.

\begin{lstlisting}[caption={Rozne operacje modyfikacji map (C++17); kompletny program.}]
#include <cassert>
#include <map>
#include <memory>
#include <string>

struct Heavy {
    static inline int constructions = 0; // C++17
    int value;
    explicit Heavy(int x) : value(x) { ++constructions; }
};

int main() {
    std::map<std::string, Heavy> m;
    auto [it, added] = m.try_emplace("a", 3);
    assert(added && it->second.value == 3);

    auto [again, added_again] = m.try_emplace("a", 4);
    assert(!added_again && again->second.value == 3);
    assert(Heavy::constructions == 1);

    m.insert_or_assign("a", Heavy{9});
    assert(m.at("a").value == 9);

    std::map<int, std::unique_ptr<int>> owned;
    owned.try_emplace(1, std::make_unique<int>(10));
    assert(*owned.at(1) == 10);
}
\end{lstlisting}

\subsection{Własny komparator i ścisły słaby porządek}

Porównanie musi spełniać warunek \emph{strict weak ordering}; w
szczególności komparator powinien być irrefleksyjny
($\neg\operatorname{comp}(x,x)$) i przechodni. Jeżeli dwie wartości
$ a,b $ nie spełniają ani \texttt{comp(a,b)}, ani \texttt{comp(b,a)},
kontener uznaje je za równoważne jako klucze --- nawet gdy
\texttt{operator==} zwraca \texttt{false}. Wymagania tego porządku
są istotne również dla \texttt{sort} i \texttt{lower\_bound}.
Na przykład komparator ignorujący wielkość liter uzna \texttt{"ABC"}
i \texttt{"abc"} za ten sam klucz: druga próba wstawienia do
\texttt{map} nie utworzy nowego elementu.

\begin{lstlisting}[caption={Mapowanie bez rozrozniania wielkosci liter: prosty komparator; kompletny program.}]
#include <cassert>
#include <cctype>
#include <map>
#include <string>

struct CaseInsensitiveLess {
    bool operator()(const std::string& a,
                    const std::string& b) const {
        auto ai = a.begin();
        auto bi = b.begin();
        while (ai != a.end() && bi != b.end()) {
            auto x = static_cast<unsigned char>(*ai++);
            auto y = static_cast<unsigned char>(*bi++);
            int lx = std::tolower(x);
            int ly = std::tolower(y);
            if (lx != ly) return lx < ly;
        }
        return a.size() < b.size();
    }
};

int main() {
    std::map<std::string, int, CaseInsensitiveLess> m;
    m.emplace("ABC", 1);
    m.emplace("abc", 2); // rownowazny klucz
    assert(m.size() == 1 && m.begin()->second == 1);
}
\end{lstlisting}

Przykład porównuje bajty w podstawowym środowisku znakowym;
\texttt{std::tolower} z bieżącą lokalizacją C nie realizuje ogólnego
Unicode case folding. Do pełnego porównywania napisów w Unicode trzeba
zastosować osobne, świadomie wybrane reguły normalizacji i kolacji.

\subsection{Wyszukiwanie heterogeniczne}

Jeżeli komparator mapy jest \emph{transparentny} (np. \texttt{std::less<>}),
możemy szukać w mapie o kluczach \texttt{std::string} za pomocą
\texttt{std::string\_view}, bez tworzenia tymczasowego ciągu znaków.
Heterogeniczne przeciążenia \texttt{find} są dostępne
od C++14, a \texttt{contains} od C++20.

\begin{lstlisting}[caption={Heterogeniczne wyszukiwanie w map (C++20); kompletny program.}]
#include <cassert>
#include <functional>
#include <map>
#include <string>
#include <string_view>

int main() {
    std::map<std::string, int, std::less<>> m{{"alpha", 3}};
    std::string_view key = "alpha";
    auto it = m.find(key);
    assert(it != m.end() && it->second == 3);
    assert(m.contains(key));
}
\end{lstlisting}

\subsection{Zmiana klucza przez uchwyt węzła (C++17)}

Kluczy w \texttt{map} nie wolno modyfikować bezpośrednio przez iterator.
Uchwyt \texttt{node\_handle} uzyskany przez \texttt{extract()} pozwala
wyjąć węzeł, zmienić jego klucz i wstawić z powrotem. Trzeba pamiętać,
że ponowne wstawienie może się nie udać z powodu kolizji kluczy.

\begin{lstlisting}[caption={Zmiana klucza przez extract (C++17); kompletny program.}]
#include <cassert>
#include <map>
#include <string>
#include <utility>

int main() {
    std::map<std::string, int> m{{"old", 7}};
    auto node = m.extract("old");
    assert(!node.empty());
    node.key() = "new";
    auto result = m.insert(std::move(node));
    assert(result.inserted && m.at("new") == 7);
}
\end{lstlisting}

\section{Kontenery haszujące: \texttt{unordered\_map} i \texttt{unordered\_set}}
\label{sec:cpp13-unordered}

Kontener \texttt{unordered\_map<Key,T,Hash,Eq>} rozmieszcza elementy
w kubełkach. Funkcja \texttt{Hash} przekształca klucz w wartość typu
\texttt{std::size\_t}; \texttt{Eq} rozstrzyga równość kluczy.
Wymagany jest warunek zgodności:
\[
\operatorname{Eq}(a,b) \Longrightarrow
\operatorname{Hash}(a)=\operatorname{Hash}(b).
\]
Warunek odwrotny nie zachodzi: kolizje haszy są dozwolone i
rozstrzygane przez porównanie kluczy. Iteracja po kontenerze
nie ma stabilnego porządku i nie należy z niego korzystać jako z
interfejsu logicznego programu.
Można wyobrazić sobie hasz jako wskazówkę, w którym kubełku szukać.
Jeśli dwa różne klucze trafią do tego samego kubełka, \texttt{Eq}
sprawdza, czy są naprawdę równe. Stała w poniższym \texttt{PointHash}
służy jedynie do wymieszania haszy obu współrzędnych.

\begin{lstlisting}[caption={Wlasny hash i porownanie kluczy (C++20); kompletny program.}]
#include <cassert>
#include <cstddef>
#include <functional>
#include <unordered_map>

struct Point {
    int x, y;
    friend bool operator==(const Point&, const Point&) = default;
};

struct PointHash {
    std::size_t operator()(const Point& p) const noexcept {
        std::size_t h1 = std::hash<int>{}(p.x);
        std::size_t h2 = std::hash<int>{}(p.y);
        return h1 ^ (h2 + std::size_t{0x9e3779b9}
                        + (h1 << 6) + (h1 >> 2));
    }
};

int main() {
    std::unordered_map<Point, int, PointHash> m;
    m[{2, 5}] = 42;
    assert(m.at(Point{2, 5}) == 42);
}
\end{lstlisting}

Ta funkcja mieszająca służy demonstracji składni, a nie zastosowaniom
kryptograficznym. Dla publicznie sterowanych kluczy przy projektowaniu
systemów wystawionych na ataki kolizyjne należy rozważyć również
odporność na złośliwe dane wejściowe.

\subsection{\texttt{load\_factor}, \texttt{reserve} i \texttt{rehash}}

\texttt{load\_factor()} to stosunek liczby elementów do liczby kubełków.
\texttt{max\_load\_factor()} podaje lub ustawia pożądany maksymalny
współczynnik. \texttt{reserve(n)} żąda przygotowania na $n$ elementów,
a \texttt{rehash(k)} na przynajmniej $k$ kubełków, z uwzględnieniem
limitów wynikających z współczynnika wypełnienia.

\begin{lstlisting}[caption={Rezerwowanie kubelek przed licznymi wstawkami (C++20); kompletny program.}]
#include <cassert>
#include <unordered_set>

int main() {
    std::unordered_set<int> seen;
    seen.reserve(1000);
    for (int i = 0; i < 1000; ++i) seen.insert(i);
    assert(seen.size() == 1000);
    assert(seen.contains(17));
}
\end{lstlisting}

Rehashing unieważnia iteratory, lecz \emph{nie unieważnia} wskaźników
i referencji do elementów. Wstawienie może unieważnić iteratory, jeśli
wywoła rehash. Usunięcie unieważnia iteratory i referencje do usuniętych
elementów. Nie wolno modyfikować kluczy w sposób zaburzający ich hash
lub równość podczas przebywania elementów w kontenerze.

\section{Adaptery: \texttt{stack}, \texttt{queue}, \texttt{priority\_queue}}
\label{sec:cpp13-adaptors}

\texttt{stack} udostępnia schemat LIFO, \texttt{queue} --- FIFO,
a \texttt{priority\_queue} --- dostęp do elementu o najwyższym
priorytecie według komparatora (domyślnie do największego).
Adaptery zwykle nie zapewniają bezpośredniej iteracji po elementach.

\begin{lstlisting}[caption={Trzy adaptery standardowe (C++11); kompletny program.}]
#include <cassert>
#include <functional>
#include <queue>
#include <stack>
#include <vector>

int main() {
    std::stack<int> s;
    s.push(1); s.push(2);
    assert(s.top() == 2);
    s.pop();

    std::queue<int> q;
    q.push(1); q.push(2);
    assert(q.front() == 1);
    q.pop();

    std::priority_queue<int, std::vector<int>, std::greater<int>> pq;
    pq.push(7); pq.push(2); pq.push(5);
    assert(pq.top() == 2); // kopiec minimum
}
\end{lstlisting}

\texttt{priority\_queue} zwykle opiera się na kopcu: \texttt{top()} ma
koszt $O(1)$, a \texttt{push()} i \texttt{pop()} $O(\log n)$. Nie jest
to kolejka FIFO ani struktura z aktualizacją priorytetu dowolnego elementu
przez gotowy publiczny interfejs.

\section{Iteratory, stabilność adresów i unieważnianie}
\label{sec:cpp13-invalidation}

Iterator nie jest ogólnie tożsamy ze wskaźnikiem. Może przechowywać
stan niezbędny do poruszania się po drzewie lub kubełkach. Po operacji
modyfikującej trzeba wiedzieć, czy iterator, referencja i wskaźnik do
elementu pozostają ważne. Typowe reguły podsumowano w tabeli.

\begin{center}
\footnotesize
\begin{tabularx}{\textwidth}{@{}lXX@{}}
\toprule
\textbf{Kontener} & \textbf{Wstawienie} & \textbf{Usunięcie} \\
\midrule
\texttt{vector} & Przy realokacji: wszystkie iteratory i referencje; bez realokacji: pozycja wstawienia i dalej & Pozycja usunięcia i dalej (także dawny \texttt{end()}) \\
\texttt{deque} & Na końcach: iteratory tracą ważność, referencje do elementów zwykle pozostają; w środku: także referencje & Zależy od pozycji; usunięcie w środku unieważnia wszystkie iteratory i referencje \\
\texttt{list} & Istniejące iteratory i referencje zachowane & Tylko usunięte elementy \\
\texttt{map}, \texttt{set} & Istniejące iteratory i referencje zachowane & Tylko usunięte elementy \\
\texttt{unordered\_*} & Iteratory mogą utracić ważność przy rehash; referencje zachowane & Tylko usunięte elementy \\
\bottomrule
\end{tabularx}
\end{center}

Dla \texttt{deque} szczególnie ostrożnie trzeba traktować
\texttt{end()}, którego ważność zależy od rodzaju usunięcia.
W tabeli podano reguły najważniejsze dla praktyki, a nie kompletną
specyfikację każdej przeciążonej operacji. Po operacjach zmieniających
strukturę bezpieczną praktyką jest korzystanie z iteratora zwróconego
przez metodę (np. \texttt{erase}, \texttt{insert}).

\subsection{Bezpieczne usuwanie podczas iteracji}

\begin{lstlisting}[caption={Usuwanie elementow map podczas iteracji (C++20); kompletny program.}]
#include <cassert>
#include <map>

int main() {
    std::map<int, int> m{{1, 2}, {2, 3}, {3, 4}};
    for (auto it = m.begin(); it != m.end();) {
        if (it->second % 2 == 0) it = m.erase(it);
        else ++it;
    }
    assert(m.size() == 1 && m.contains(2));
}
\end{lstlisting}

\paragraph{Szczególny przypadek \texttt{vector<bool>}.}
\texttt{std::vector<bool>} jest specjalizacją oszczędzającą pamięć,
która może reprezentować wartości w postaci bitów. Jej \texttt{operator[]}
nie musi zwracać \texttt{bool\&}, tylko specjalny obiekt proxy. Nie
należy oczekiwać normalnej tablicy obiektów \texttt{bool}, ciągłego
przechowywania pojedynczych \texttt{bool} ani bezpiecznej modyfikacji
odrębnych bitów z różnych wątków bez synchronizacji.

\section{Algorytmy standardowe i wymagania wobec iteratorów}

\begin{center}\small
\begin{tabularx}{\textwidth}{@{}p{0.34\textwidth}X@{}}
\toprule
Algorytm & Działanie i warunki \\
\midrule
\code{std::find} & Szuka pierwszej wartości równej wzorcowi; przy braku zwraca iterator końcowy. \\
\code{std::transform} & Stosuje funkcję do zakresu i zapisuje wyniki przez iterator wyjściowy. \\
\code{std::sort} & Sortuje zakres z iteratorami dostępu swobodnego; nie zachowuje kolejności elementów równoważnych. \\
\code{std::stable_sort} & Sortuje, zachowując wzajemną kolejność elementów równoważnych. \\
\code{std::lower_bound} & Szuka pierwszego elementu nie mniejszego od wzorca; wymaga odpowiednio uporządkowanego zakresu. \\
\code{std::erase_if} & Usuwa z kontenera elementy spełniające predykat; dostępne od C++20 dla wspieranych kontenerów. \\
\bottomrule
\end{tabularx}
\end{center}
\label{sec:cpp13-algorithms}

Algorytmy oddzielają \emph{operację} od \emph{reprezentacji danych}.
Wersje klasyczne przyjmują iteratory, a w C++20 wiele algorytmów
w \texttt{std::ranges} można wywołać bezpośrednio na zakresie. Szczegóły
kategorii iteratorów i widoków omawiamy w rozdziale~\ref{chap:ranges-views-cpp}; tutaj skupiamy
się na doborze algorytmu i jego kontrakcie.

\subsection{Wyszukiwanie, transformacje i agregacje}

Do częstych operacji należą \texttt{find}, \texttt{find\_if},
\texttt{count\_if}, \texttt{transform}, \texttt{accumulate},
\texttt{all\_of}, \texttt{any\_of}, \texttt{none\_of} i
\texttt{for\_each}. Algorytmy numeryczne znajdują się w nagłówku
\texttt{<numeric>}, pozostałe głównie w \texttt{<algorithm>}.

\begin{lstlisting}[caption={Algorytmy klasyczne i Ranges (C++20); kompletny program.}]
#include <algorithm>
#include <cassert>
#include <numeric>
#include <vector>

int main() {
    std::vector<int> values{1, 2, 3, 4, 5};
    auto even = [](int x) { return x % 2 == 0; };

    assert(std::ranges::count_if(values, even) == 2);
    assert(std::ranges::all_of(values,
        [](int x) { return x > 0; }));
    assert(std::accumulate(values.begin(), values.end(), 0) == 15);

    std::vector<int> squares(values.size());
    std::ranges::transform(values, squares.begin(),
                           [](int x) { return x * x; });
    assert(squares[2] == 9);
}
\end{lstlisting}

\subsection{Sortowanie i stabilność kolejności}

\texttt{std::sort} wymaga iteratorów dostępu swobodnego i ma
złożoność $O(n\log n)$ w najgorszym przypadku (w liczbie porównań).
\texttt{stable\_sort} zachowuje kolejność elementów, które komparator
uznaje za równoważne, zwykle kosztem dodatkowych zasobów; jego
szczegółowe gwarancje złożoności zależą od dostępnej pamięci.
\texttt{nth\_element} umieszcza wybrany element na pozycji, którą
zajmowałby po sortowaniu, bez sortowania całej sekwencji; przeciętna
liczba porównań jest liniowa.

\begin{lstlisting}[caption={Stabilne sortowanie rekordow po kluczu (C++20); kompletny program.}]
#include <algorithm>
#include <cassert>
#include <string>
#include <vector>

struct Record {
    std::string name;
    int group;
};

int main() {
    std::vector<Record> v{{"A", 2}, {"B", 1}, {"C", 2}};
    std::ranges::stable_sort(v, {}, &Record::group);
    assert(v[0].name == "B");
    assert(v[1].name == "A" && v[2].name == "C");
}
\end{lstlisting}

\subsection{Wyszukiwanie binarne i porządek wejścia}

\texttt{binary\_search}, \texttt{lower\_bound}, \texttt{upper\_bound}
i \texttt{equal\_range} wymagają odpowiedniego uporządkowania
(lub partycjonowania względem danego klucza) zakresu. Zastosowanie
do nieuporządkowanych danych daje wynik niespełniający kontraktu
algorytmu. W wersjach na iteratorach nie-losowego dostępu liczba
porównań może pozostać logarytmiczna, lecz przesuwanie iteratorów
może być liniowe.

\begin{lstlisting}[caption={Wyszukiwanie miejsca wstawienia w posortowanym vector (C++11); kompletny program.}]
#include <algorithm>
#include <cassert>
#include <vector>

int main() {
    std::vector<int> sorted{1, 3, 3, 7};
    auto it = std::lower_bound(sorted.begin(), sorted.end(), 3);
    assert(it == sorted.begin() + 1);

    auto upper = std::upper_bound(sorted.begin(), sorted.end(), 3);
    assert(upper == sorted.begin() + 3);
    assert(std::binary_search(sorted.begin(), sorted.end(), 7));
}
\end{lstlisting}

\subsection{Idiom \emph{erase--remove} i \texttt{std::erase\_if}}

\texttt{std::remove} nie usuwa elementów \emph{z kontenera}:
przestawia elementy pozostające i zwraca logiczny nowy koniec.
Rozmiar kontenera zmniejsza dopiero wywołanie \texttt{erase}.
Od C++20 wygodniej używać \texttt{std::erase} i
\texttt{std::erase\_if} dla wspieranych kontenerów.

\begin{lstlisting}[caption={Dwa sposoby usuwania z vector (C++20); kompletny program.}]
#include <algorithm>
#include <cassert>
#include <vector>

int main() {
    std::vector<int> a{1, 2, 3, 4, 5};
    auto last = std::remove_if(a.begin(), a.end(),
                               [](int x) { return x % 2 == 0; });
    a.erase(last, a.end());
    assert((a == std::vector<int>{1, 3, 5}));

    std::vector<int> b{1, 2, 3, 4, 5};
    auto removed = std::erase_if(b,
                                 [](int x) { return x % 2 == 0; });
    assert(removed == 2 && b == a);
}
\end{lstlisting}

Dla kontenerów asocjacyjnych nie wolno stosować klasycznego
\texttt{remove\_if} tak jak dla wektora: ich elementy mają klucze,
których pozycja jest regulowana przez strukturę danych. Stosujemy
\texttt{erase} z iteratorem lub \texttt{std::erase\_if} (C++20).

\subsection{Algorytmy kopca i \texttt{std::priority\_queue}}

\texttt{make\_heap}, \texttt{push\_heap}, \texttt{pop\_heap} i
\texttt{sort\_heap} pozwalają użyć własnego wektora jako kopca.
\texttt{make\_heap} działa w czasie liniowym, a operacje dopisania i
usunięcia maksimum w $O(\log n)$. Są to operacje inne niż pełne sortowanie.

\section{Przypadki projektowe: dobór kontenera do problemu}
\label{sec:cpp13-use-cases}

\subsection{Histogram, deduplikacja i lista wyników}

Załóżmy, że chcemy zebrać wartości, usunąć duplikaty i otrzymać
wynik w kolejności rosnącej. Możemy użyć \texttt{std::set}, albo
\texttt{vector} z \texttt{sort} i \texttt{unique}. Drugi sposób bywa
korzystniejszy w przetwarzaniu wsadowym, ponieważ \texttt{vector}
przechowuje dane ciągle i ogranicza liczbę alokacji węzłów.

\begin{lstlisting}[caption={Deduplikacja przez sort i unique (C++11); kompletny program.}]
#include <algorithm>
#include <cassert>
#include <vector>

int main() {
    std::vector<int> keys{5, 2, 5, 1, 2, 8};
    std::sort(keys.begin(), keys.end());
    keys.erase(std::unique(keys.begin(), keys.end()), keys.end());
    assert((keys == std::vector<int>{1, 2, 5, 8}));
}
\end{lstlisting}

\subsection{Własność elementów w kontenerach polimorficznych}

Przechowywanie obiektów klas pochodnych przez wartość w
\texttt{vector<Base>} może powodować \emph{object slicing}. Dla
polimorfizmu dynamicznego używamy np. wektora wskaźników
\texttt{unique\_ptr<Base>}. Przy zamkniętej rodzinie typów można użyć
wektora z elementami \texttt{variant} (zob. rozdział~\ref{chap:cpp14-type-erasure}).
Różnicę między polimorfizmem a obcięciem obiektu omawia
rozdział~\ref{chap:cpp-object-model}.

\begin{lstlisting}[caption={Kontener wlasciwych obiektow polimorficznych (C++14); kompletny program.}]
#include <cassert>
#include <memory>
#include <vector>

struct Shape {
    virtual ~Shape() = default;
    virtual int area() const = 0;
};
struct Square : Shape {
    int side;
    explicit Square(int x) : side(x) {}
    int area() const override { return side * side; }
};

int main() {
    std::vector<std::unique_ptr<Shape>> shapes;
    shapes.push_back(std::make_unique<Square>(3));
    assert(shapes.front()->area() == 9);
}
\end{lstlisting}

\subsection{Kontenery a współbieżność}

Samo wybranie kontenera standardowego nie czyni jego operacji
modyfikujących bezpiecznymi dla współbieżnego wykonania. Równoczesne
modyfikowanie wspólnego \texttt{vector} przez niezależne
\texttt{push\_back()} bez synchronizacji jest wyścigiem danych.
Można natomiast równolegle odczytywać elementy kontenera, jeśli żaden
wątek ich nie modyfikuje, a czas życia kontenera jest zachowany.
Dla wielu kontenerów można także modyfikować \emph{różne istniejące}
elementy równolegle, jeśli nie dochodzi do konfliktu pamięci; wyjątek
stanowi m.in. specjalizacja \texttt{vector<bool>}. Zmiany strukturalne
wymagają osobnej analizy i synchronizacji (zob. rozdział~\ref{chap:concurrency}).

\section{Złożoność, lokalność pamięci i pomiar wydajności}
\label{sec:cpp13-performance}

Teoretyczna złożoność to jeden z elementów oceny. Istotne są również:
liczba alokacji, liczba odwołań do rozproszonych adresów, rozmiar
elementu, przewidywalność gałęzi i zachowanie cache procesora.
Dla wielu zastosowań \texttt{vector} z wyszukiwaniem liniowym
może być szybszy od struktury drzewiastej przy małej liczbie kluczy.
To nie przeczy złożoności asymptotycznej: stałe i charakter obciążenia
są częścią kosztu rzeczywistego programu.

\subsection{Koszt sumowania ciągłych danych}

\begin{lstlisting}[caption={Prosty szkielet pomiaru vector i list (C++20); kompletny program.}]
#include <chrono>
#include <cstdint>
#include <iostream>
#include <list>
#include <numeric>
#include <utility>
#include <vector>

int main() {
    constexpr int N = 200000;
    std::vector<int> v(N, 1);
    std::list<int> l(v.begin(), v.end());

    auto measure = [](const auto& xs) {
        auto begin = std::chrono::steady_clock::now();
        std::int64_t sum = std::accumulate(xs.begin(), xs.end(),
                                           std::int64_t{0});
        auto end = std::chrono::steady_clock::now();
        auto ns = std::chrono::duration_cast<
            std::chrono::nanoseconds>(end - begin).count();
        return std::pair{sum, ns};
    };

    auto [sum_v, time_v] = measure(v);
    auto [sum_l, time_l] = measure(l);
    std::cout << "vector: " << sum_v << ", " << time_v << " ns\n";
    std::cout << "list:   " << sum_l << ", " << time_l << " ns\n";
}
\end{lstlisting}

To \emph{ilustracja sposobu pomiaru}, a nie miarodajny benchmark.
Jednorazowy wynik bywa zdominowany przez rozgrzanie cache, optymalizacje
kompilatora, scheduler i zakłócenia systemowe. Poprawniejszy eksperyment
powtarza pomiary, rozdziela przygotowanie danych od mierzonej operacji,
uwzględnia pamięć, różne rozmiary danych, medianę i rozrzut pomiarów.
Dla wiarygodnych porównań warto użyć np. Google Benchmark i profilera.

\subsection{Przybliżone reguły wyboru}

\begin{center}
\small
\begin{tabularx}{\textwidth}{@{}lX@{}}
\toprule
\textbf{Wymaganie} & \textbf{Rozsądny wybór początkowy} \\
\midrule
Dużo iteracji, obliczeń, dostęp indeksowany & \texttt{vector}; przy stałym rozmiarze \texttt{array} \\
Częste dodawanie i usuwanie z obu końców & \texttt{deque} \\
Utrzymanie kolejności kluczy, wyszukiwanie przedziałowe & \texttt{map} / \texttt{set} \\
Wyszukiwanie po kluczu bez wymogu kolejności & \texttt{unordered\_map} / \texttt{unordered\_set} \\
Przenoszenie istniejących węzłów i stabilne iteratory & \texttt{list} \\
FIFO / LIFO & \texttt{queue} / \texttt{stack} \\
Wielokrotne pobieranie największego lub najmniejszego & \texttt{priority\_queue} \\
\bottomrule
\end{tabularx}
\end{center}

\section{Najczęstsze błędy i dobre praktyki}
\label{sec:cpp13-pitfalls}

\begin{enumerate}
\item Używanie \texttt{vector::operator[]} bez pewności, że indeks jest
  poprawny; w kodzie granicznym warto stosować \texttt{at()}.
\item Traktowanie \texttt{reserve} jako tworzenia elementów (mylenie
  \texttt{size} z \texttt{capacity}).
\item Korzystanie z iteratorów lub referencji po unieważnieniu przez
  realokację, \texttt{erase} lub \texttt{rehash}.
\item Zakładanie, że \texttt{emplace\_back} zawsze jest szybsze od
  \texttt{push\_back}.
\item Używanie \texttt{map::operator[]} do samego sprawdzania istnienia
  klucza, z niezamierzonym wstawieniem nowej wartości.
\item Pisanie komparatorów naruszających \emph{strict weak ordering};
  komparator \texttt{a <= b} nie jest poprawną relacją porządkującą.
\item Używanie \texttt{std::sort} na \texttt{std::list} lub
  \texttt{std::lower\_bound} na danych bez odpowiedniego porządku.
\item Wnioskowanie o kolejności elementów \texttt{unordered\_map}
  na podstawie przypadkowego przebiegu programu.
\item Niepotrzebne stosowanie \texttt{list} z przekonaniem, że
  wstawienie $O(1)$ zawsze wygrywa z \texttt{vector}.
\item Wprowadzanie równoczesnych modyfikacji kontenera z wielu wątków
  bez uzasadnienia braku konfliktu i odpowiedniej synchronizacji.
\end{enumerate}

\section{Pytania kontrolne}

\begin{enumerate}
\item Jaka jest różnica między \texttt{reserve} a \texttt{resize}, i
  które z nich tworzy nowe elementy?
\item Kiedy \texttt{push\_back} na \texttt{vector} unieważnia wszystkie
  referencje, a kiedy zachowuje referencje do wcześniejszych elementów?
\item Dlaczego \texttt{map::operator[]} może być niepożądany podczas
  sprawdzania istnienia klucza?
\item Jakie wymagania muszą spełniać \texttt{Hash} i \texttt{Eq} w
  \texttt{unordered\_map}?
\item Dlaczego \texttt{std::remove} nie zmienia rozmiaru \texttt{vector}?
\item Jak różnią się koszty i kontrakty \texttt{sort}, \texttt{stable\_sort},
  \texttt{nth\_element} i \texttt{lower\_bound}?
\item W jakich okolicznościach \texttt{list::splice} zapewnia korzyść
  niedostępną dla zwykłego \texttt{vector}?
\end{enumerate}

\section{Podsumowanie rozdziału}

Dobry kod C++ wybiera kontener zgodnie z potrzebnymi operacjami i
własnościami pamięci, a nie jedynie na podstawie pojedynczego rzędu
złożoności. \texttt{vector} to często najlepszy punkt wyjścia;
\texttt{deque}, kontenery drzewiaste, haszujące i listy rozwiązują
bardziej szczegółowe problemy. Praktycznie równie ważne jak złożoności
są reguły unieważniania iteratorów, wymagania komparatorów,
bezpieczeństwo czasu życia i świadome wykorzystywanie algorytmów
standardowych zamiast ręcznie pisanych pętli.

\section*{Literatura i dokumentacja}
\addcontentsline{toc}{section}{Literatura i dokumentacja}
\begin{itemize}
\item \url{https://en.cppreference.com/w/cpp/container} --- kontenery i
  zestawienie reguł unieważniania iteratorów.
\item \url{https://en.cppreference.com/w/cpp/algorithm} --- algorytmy
  biblioteki standardowej i ich kontrakty.
\item \url{https://en.cppreference.com/w/cpp/container/vector} ---
  \texttt{std::vector} i zarządzanie pojemnością.
\item \url{https://en.cppreference.com/w/cpp/container/unordered_map} ---
  struktury haszujące.
\item Scott Meyers, \emph{Effective Modern C++}, O'Reilly, 2014.
\item Nicolai M. Josuttis, \emph{The C++ Standard Library}, wydanie drugie,
  Addison-Wesley, 2012.
\end{itemize}
